% Minimal preamble for Open Logic Project sources (pinned commit 1e960beff9ed7835bf3e3f1335e21af3439cd107).
%
% OLP's own master preamble cannot run on the project host: TeX Live 2025 there lacks
% prooftree.sty, ntheorem.sty, extarrows.sty, mathdots.sty, isomath.sty and subfiles.sty.
% Following the project's own rule ("if the author preamble needs packages the toolchain lacks,
% do not invent macros -- use a minimal working preamble"), this file declares, verbatim in
% spirit and numbering behaviour, exactly the theorem-like environments OLP defines in
% sty/open-logic-book-envs.sty (thm/lem/prop/cor/defn/ex/prob/rem/note/case/axiom/conv),
% plus the shorthand accents OLP uses (`!!a`, `!!b`, `!!{...}`).
%
% Author text is untouched; only the packaging is ours.
\documentclass[10pt,a4paper]{article}

\usepackage{amsmath}
\usepackage{amssymb}
\usepackage{amsthm}
\usepackage{graphicx}
\usepackage{hyperref}
\usepackage{enumitem}
\usepackage[framemethod=default]{mdframed}

% ---- OLP shorthand accents (from sty/open-logic.sty) ----
% NB: \! is taken by LaTeX (negative thin space), so we only define the !!a / !!b families
% and the \olup helper used by them.
\newcommand*{\olup}[1]{\MakeUppercase{#1}}
\def\!!!a#1{\olup{#1}}
\def\!a#1{\olup{#1}}
\def\!b#1{\olup{#1}}

% ---- theorem-like environments, same family numbering as OLP ----
\theoremstyle{plain}
\newtheorem{thm}{Theorem}[section]
\newtheorem{lem}[thm]{Lemma}
\newtheorem{prop}[thm]{Proposition}
\newtheorem{cor}[thm]{Corollary}
\newtheorem{axiom}{Axiom}

\theoremstyle{definition}
\newtheorem{defn}[thm]{Definition}
\newtheorem{ex}[thm]{Example}
\newtheorem{prob}{Problem}[section]
\newtheorem{probd}{Problem}
\newtheorem{conv}[thm]{Convention}

\theoremstyle{remark}
\newtheorem{rem}[thm]{Remark}
\newtheorem{note}[thm]{Note}
\newtheorem{case}[thm]{Case}

% ---- environments OLP defines as wrappers ----
\newenvironment{probtag}[1]{\begin{prob}}{\end{prob}}
\newenvironment{explain}{}{}
\newenvironment{explain*}{}{}
\newenvironment{editorial}{}{}
\newenvironment{digress}{}{}
\newenvironment{digress*}{}{}
\newenvironment{history}{}{}
\newenvironment{reference}{}{}
\newcommand*{\olfileid}[3]{}
\newcommand*{\olsection}[1]{\section*{#1}}
\newcommand*{\olchapter}[3]{}
\newcommand*{\olpart}[2]{}
\newcommand*{\olimport}[1]{}
\newcommand*{\olinclude}[1]{}
\newcommand*{\olpath}[1]{}
\newcommand*{\ollangid}[1]{}
\newcommand*{\problemtag}[1]{}
\newcommand*{\problemauthor}[1]{}
\newcommand*{\ollabel}[1]{\label{#1}}
% \olref takes 0-3 OPTIONAL arguments plus one mandatory label:
%     \olref{lab}  |  \olref[from]{lab}  |  \olref[from][to]{lab}  |  \olref[a][b][c]{lab}
% A single-argument definition silently swallows "[sth]" and emits a broken reference
% (measured on the host: the text came out as "??iso]x"). The chain below consumes the
% brackets and cross-references only the mandatory argument.
\makeatletter
\newcommand{\olref}{\@ifnextchar[{\olref@i}{\olref@final}}
\def\olref@i[#1]{\@ifnextchar[{\olref@ii}{\olref@final}}
\def\olref@ii[#1]{\@ifnextchar[{\olref@iii}{\olref@final}}
\def\olref@iii[#1]{\olref@final}
\def\olref@final#1{\ref{#1}}
\makeatother
\newcommand*{\olnewpage}{\clearpage}

% ===================== OLP compatibility layer (1) =====================

\makeatletter

% ---------------- proof trees (OLP's prooftree is a bussproofs derivative) ----------------
\providecommand{\AxiomC}[1]{}
\providecommand{\UnaryInfC}[1]{}
\providecommand{\BinaryInfC}[1]{}
\providecommand{\TrinaryInfC}[1]{}
\providecommand{\QuaternaryInfC}[1]{}
\providecommand{\QuinaryInfC}[1]{}
\providecommand{\RightLabel}[1]{}
\providecommand{\LeftLabel}[1]{}
\providecommand{\doubleLine}{}
\providecommand{\noLine}{}
\providecommand{\DisplayProof}{}
\providecommand{\fCenter}{}
\newenvironment{prooftree}{}{}
\newenvironment{ltprooftree}{}{}
\newenvironment{oltableau}{}{}

% ---------------- logical notation ----------------
\providecommand{\lif}{\mathbin{\rightarrow}}
\providecommand{\lforall}{\mathord{\forall}}
\providecommand{\lexists}{\mathord{\exists}}
\providecommand{\ltrue}{\top}
\providecommand{\lfalse}{\bot}
\providecommand{\Ltrue}{\top}
\providecommand{\Lfalse}{\bot}
\providecommand{\Sequent}{\mathrel{\Rightarrow}}
\providecommand{\Proves}{\mathrel{\vdash}}
\providecommand{\Entails}{\mathrel{\models}}
\providecommand{\Subst}[3]{#1[#2/#3]}
\providecommand{\sFmla}[2]{#1(#2)}
\providecommand{\pFmla}[2]{#1(#2)}
\providecommand{\formula}{\relax}
\providecommand{\True}{\top}
\providecommand{\False}{\bot}
\providecommand{\Args}[1]{#1}
\providecommand{\Ts}{\ensuremath{\top}}
\providecommand{\Fi}{\ensuremath{\bot}}
\providecommand{\lfalsum}{\bot}

% ---------------- objects, structures, metatheory ----------------
\providecommand{\Obj}{}
\providecommand{\Struct}{\ensuremath{\mathfrak{A}}}
\providecommand{\Domain}{\ensuremath{|\mathfrak{A}|}}
\providecommand{\Assign}{\ensuremath{s}}
\providecommand{\Sat}[2]{\ensuremath{#1 \models #2}}
\providecommand{\mSat}[3]{\ensuremath{#1 \models #2[#3]}}
\providecommand{\mModel}[2]{\ensuremath{#1 \models #2}}
\providecommand{\Lang}{\ensuremath{\mathcal{L}}}
\providecommand{\Th}{\ensuremath{\mathrm{Th}}}
\providecommand{\Log}{\ensuremath{\mathrm{L}}}
\providecommand{\tuple}[1]{\langle #1 \rangle}
\providecommand{\num}[1]{\underline{#1}}
\providecommand{\cfind}[1]{\varphi_{#1}}
\providecommand{\Nat}{\ensuremath{\mathbb{N}}}
\providecommand{\Mult}{\ensuremath{\mathrm{Mult}}}
\providecommand{\Add}{\ensuremath{\mathrm{Add}}}
\providecommand{\Exp}{\ensuremath{\mathrm{Exp}}}
\providecommand{\Zero}{\ensuremath{0}}
\providecommand{\Succ}{\ensuremath{\mathrm{S}}}
\providecommand{\lambd}{\ensuremath{\lambda}}
\providecommand{\bredpar}{\mathrel{\rightarrow_\beta}}
\providecommand{\bcd}[1]{\ensuremath{\ulcorner #1 \urcorner}}
\providecommand{\Set}[1]{\ensuremath{\{#1\}}}
\providecommand{\indcase}[1]{\par\noindent\emph{Case #1}.\quad}
\providecommand{\indfrm}[1]{\par\noindent\emph{Form #1}.\quad}
\providecommand{\iftag}[1]{#1}
\providecommand{\tagitem}[1]{\item[#1]}
\providecommand{\Elim}[1]{}
\providecommand{\Intro}[1]{}
\providecommand{\LeftR}[1]{}
\providecommand{\RightR}[1]{}
\providecommand{\TAss}[1]{}
\providecommand{\TRule}[1]{}

\makeatother
\providecommand{\olpath}{.}

% ===================== OLP compatibility layer (2) =====================
% --- second compatibility block: notation and shorthands used by the frozen OLP units ---
\providecommand{\CutCS}{\mathrm{Cut}}
\providecommand{\LeftSubproofLabel}{\relax}
\providecommand{\RightSubproofLabel}{\relax}
\providecommand{\maeh}{\relax}
\providecommand{\cheight}{\relax}
\providecommand{\ande}{\mathrm{and}_E}
\providecommand{\andi}{\mathrm{and}_I}
\providecommand{\ore}{\mathrm{or}_E}
\providecommand{\ori}{\mathrm{or}_I}
\providecommand{\typeof}{\mathord{:}}
\providecommand{\pValue}[1]{#1}
\providecommand{\Sent}{\mathrm{Sent}}
\providecommand{\cardplus}{\mathbin{+}}
\providecommand{\cardtimes}{\mathbin{\cdot}}
\providecommand{\cardexpo}{\mathbin{\hat{\ }}}
\providecommand{\cardnless}{\not<}
\providecommand{\cardle}{\leq}
\providecommand{\funfromto}[3]{#1\colon #2 \to #3}
\providecommand{\funimage}[1]{\operatorname{im}(#1)}
\providecommand{\fdefined}{\mathord{\downarrow}}
\providecommand{\canonord}{\operatorname{ord}}
\providecommand{\ordtimes}{\mathbin{\cdot}}
\providecommand{\ordeq}[2]{#1 = #2}
\providecommand{\ordneq}[2]{#1 \neq #2}
\providecommand{\ordtype}[1]{\operatorname{ord}(#1)}
\providecommand{\ordsucc}{\mathord{+1}}
\providecommand{\ordunion}{\cup}
\providecommand{\onesphere}{S^1}
\providecommand{\rotationsgroup}{\mathrm{SO}(2)}
\providecommand{\equivrep}[1]{[#1]}
\providecommand{\rlexless}{\prec}
\providecommand{\ZFC}{\mathrm{ZFC}}
\providecommand{\ZF}{\mathrm{ZF}}
\providecommand{\Rat}{\mathbb{Q}}
\providecommand{\Real}{\mathbb{R}}
\providecommand{\PosInt}{\mathbb{Z}^{+}}
\providecommand{\Id}{\mathrm{id}}
\providecommand{\Bin}{\mathrm{Bin}}
\providecommand{\pto}{\rightharpoonup}
\providecommand{\dom}{\operatorname{dom}}
\providecommand{\ran}{\operatorname{ran}}
\providecommand{\len}{\operatorname{len}}
\providecommand{\concat}{\mathbin{{}^\frown}}
\providecommand{\card}[1]{\lvert #1 \rvert}
\providecommand{\Pow}{\mathcal{P}}
\providecommand{\Setabs}[2]{\{\, #1 \mid #2 \,\}}
\providecommand{\comp}{\circ}
\providecommand{\Char}{\chi}
\providecommand{\Prf}{\operatorname{Prf}}
\providecommand{\Prov}{\operatorname{Prov}}
\providecommand{\Proj}{\operatorname{proj}}
\providecommand{\Gn}{\operatorname{gn}}
\providecommand{\gn}{\operatorname{gn}}
\providecommand{\Value}[1]{\operatorname{val}(#1)}
\providecommand{\Deduce}{\mathrm{D}}
\providecommand{\Dual}{\mathrm{d}}
\providecommand{\Ax}{\mathrm{Ax}}
\providecommand{\pSat}[3]{#1 \models #2\,[#3]}
\providecommand{\pAssign}{\sigma}
\providecommand{\bforall}{\forall}
\providecommand{\bexists}{\exists}
\providecommand{\liff}{\mathbin{\leftrightarrow}}
\providecommand{\FV}{\operatorname{FV}}
\providecommand{\UnaryInf}{\ensuremath{\mid}}
\providecommand{\BinaryInf}{\ensuremath{\mid}}
\providecommand{\depth}{\operatorname{depth}}
\providecommand{\oliflabeldef}[3]{#1}
\providecommand{\TMblank}{\sqcup}
\providecommand{\TMstroke}{\mid}
\providecommand{\TMstay}{\mathrm{S}}
\providecommand{\TMright}{\mathrm{R}}
\providecommand{\TMleft}{\mathrm{L}}
\providecommand{\TMendtape}{\triangleright}
\providecommand{\umin}{\mu}
\providecommand{\Sum}{\sum}
\providecommand{\citep}[1]{\cite{#1}}
\providecommand{\citet}[1]{\cite{#1}}
\providecommand{\citeauthor}[1]{\cite{#1}}
\providecommand{\citeyear}[1]{\cite{#1}}

% ===================== OLP compatibility layer (3) =====================
% --- third compatibility block: remaining notation seen in the frozen OLP units ---
\providecommand{\iso}{\cong}
\providecommand{\nszero}{\mathrm{0}}
\providecommand{\nssucc}{\mathrm{S}}
\providecommand{\nsplus}{\mathbin{+}}
\providecommand{\nssless}{<}
\providecommand{\nsless}{<}
\providecommand{\PL}{\mathrm{PL}}
\providecommand{\Nec}{\mathord{\Box}}
\providecommand{\RK}{\mathrm{RK}}
\providecommand{\ST}{\mathrm{ST}}
\providecommand{\SSubst}[3]{#1[#2/#3]}
\providecommand{\subst}[3]{#1[#2/#3]}
\providecommand{\tagfalse}{\mathrm{F}}
\providecommand{\cutr}{\mathrm{Cut}}
\providecommand{\cardneq}{\neq}
\providecommand{\issucc}[1]{\operatorname{succ}(#1)}
\providecommand{\tup}[1]{\langle #1 \rangle}
\providecommand{\Var}{\mathrm{Var}}
\providecommand{\Bew}{\operatorname{Bew}}
\providecommand{\Con}{\operatorname{Con}}
\providecommand{\Cons}{\operatorname{Cons}}
\providecommand{\pDeduce}{\mathrm{D}}
\providecommand{\psat}{\models}
\providecommand{\den}[1]{\llbracket #1 \rrbracket}
\providecommand{\Sem}[2]{\llbracket #1 \rrbracket_{#2}}
\providecommand{\DomainOf}[1]{\lvert #1 \rvert}
\providecommand{\Form}{\mathrm{Form}}
\providecommand{\Terms}{\mathrm{Term}}
\providecommand{\Atom}{\mathrm{Atom}}
\providecommand{\FVars}[1]{\operatorname{FV}(#1)}
\providecommand{\fv}[1]{\operatorname{fv}(#1)}
\providecommand{\tuplep}[1]{\langle #1 \rangle}
\providecommand{\restrict}[2]{#1 \upharpoonright #2}
\providecommand{\sto}{\rightarrow}
\providecommand{\FmlaSet}[1]{\mathrm{Fmla}_{#1}}
\providecommand{\Thm}{\mathrm{Thm}}
\providecommand{\Eval}{\operatorname{val}}
\providecommand{\blank}{\sqcup}
\providecommand{\blankTape}{\sqcup}
\providecommand{\Consis}{\operatorname{Con}}
\providecommand{\drs}{\mathrm{d}}
\providecommand{\VRule}{\mid}
\providecommand{\Bet}{\mathrm{B}}
\providecommand{\Seg}[1]{\mathrm{Seg}_{#1}}
\providecommand{\Enums}{\mathrm{En}}
\providecommand{\Z}{\mathbb{Z}}
\providecommand{\N}{\mathbb{N}}
\providecommand{\abs}[1]{\lvert #1 \rvert}
\providecommand{\floor}[1]{\lfloor #1 \rfloor}
\providecommand{\ceil}[1]{\lceil #1 \rceil}

% ===================== OLP compatibility layer (4) =====================
% --- fourth compatibility block: harvested from the per-candidate compile logs ---
\providecommand{\beredpar}{\mathrel{\rightarrow_\beta}}
\makeatletter
\def\bered{\@ifnextchar[{\bered@i}{\mathrel{\rightarrow_\beta}}}
\def\bered@i[#1]{\mathrel{\rightarrow_\beta}}
\def\bred{\@ifnextchar[{\bred@i}{\mathrel{\rightarrow_\beta}}}
\def\bred@i[#1]{\mathrel{\rightarrow_\beta}}
\makeatother
\providecommand{\beredone}{\mathrel{=_\beta}}
\providecommand{\bredone}{\mathrel{=_\beta}}
\providecommand{\xredone}{\mathrel{=_\beta}}
\providecommand{\aconvone}{\mathrel{=_\alpha}}
\providecommand{\cardless}{<}
\providecommand{\tagrefs}[1]{#1}
\providecommand{\ordplus}{\mathbin{+}}
\providecommand{\funrestrictionto}[2]{#1 \upharpoonright #2}
\providecommand{\fundefined}{\mathord{\uparrow}}
\providecommand{\Trm}{\mathrm{Trm}}
\providecommand{\OProv}{\operatorname{Prov}}
\providecommand{\OPrf}{\operatorname{Prf}}
\providecommand{\aeq}{\mathrel{=}}
\providecommand{\mClass}[1]{[#1]}
\providecommand{\equivclass}[1]{[#1]}
\providecommand{\pheight}{\operatorname{ht}}
\providecommand{\phy}{\operatorname{ht}}
\providecommand{\Cut}{\mathrm{Cut}}
\providecommand{\QuantRank}{\operatorname{qr}}
\providecommand{\elemequiv}{\mathrel{\equiv}}
\providecommand{\PIso}{\mathrm{PIso}}
\providecommand{\Theory}{\mathrm{Th}}
\providecommand{\supstrict}{\subsetneq}
\providecommand{\trcl}{\operatorname{trcl}}
\providecommand{\cref}[1]{\ref{#1}}
\providecommand{\Discharge}{\mathrm{Dis}}
\providecommand{\Weakening}{\mathrm{Wk}}
\providecommand{\ZFminus}{\mathrm{ZF}^{-}}
\makeatletter
\renewcommand{\Dual}{\@ifnextchar[{\Dual@i}{\ifmmode\mathrm{d}\else$\mathrm{d}$\fi}}
\def\Dual@i[#1]{\ifmmode\mathrm{d}[#1]\else$\mathrm{d}[#1]$\fi}
\makeatother
\providecommand{\RightR}{\mathrm{R}}
\providecommand{\LeftR}{\mathrm{L}}
\providecommand{\Intro}{\mathrm{I}}
\providecommand{\Axiom}{\mathrm{Ax}}
\providecommand{\formula}[1]{#1}
\providecommand{\stared}{\ast}
\providecommand{\oset}[2]{\{\, #1 \mid #2 \,\}}
\providecommand{\setof}[2]{\{\, #1 \mid #2 \,\}}
\providecommand{\pow}[1]{\mathcal{P}(#1)}
\providecommand{\ordpair}[2]{\langle #1, #2 \rangle}
\providecommand{\ordtriple}[3]{\langle #1, #2, #3 \rangle}
\providecommand{\union}{\cup}
\providecommand{\inter}{\cap}
\providecommand{\seq}[1]{\langle #1 \rangle}
\providecommand{\SuchThat}{\mid}
\providecommand{\where}{\mid}
\providecommand{\TVar}{\mathrm{TVar}}
\providecommand{\Bool}{\mathrm{Bool}}
\providecommand{\Val}{\operatorname{Val}}
\providecommand{\TL}[1]{\mathrm{TL}}
\providecommand{\xstared}{\ast}

% ===================== OLP compatibility layer (5) =====================
% --- fifth compatibility block: rescues the quarantined OLP candidates ---
\providecommand{\ident}{\operatorname{id}}
\providecommand{\setrank}{\operatorname{rank}}
\providecommand{\cardfont}{\mathfrak{c}}
\providecommand{\Complement}{\overline}
\providecommand{\QR}{\mathrm{QR}}
\providecommand{\becd}{\ensuremath{\ulcorner}}
\providecommand{\disjointsum}{\mathbin{\oplus}}
\providecommand{\cardeq}[2]{#1 = #2}
\providecommand{\xred}{\mathrel{\rightarrow}}
% --- sixth compatibility block: names reported by real compile logs ---
\providecommand{\ORefut}{\operatorname{Refut}}
\providecommand{\ORProv}{\operatorname{Prov}}
\providecommand{\cutrank}{\operatorname{cr}}
\providecommand{\Refut}{\operatorname{Refut}}
\providecommand{\DischargeRule}{\mathrm{Dis}}
\providecommand{\Contraction}{\mathrm{C}}
\providecommand{\bottomAlignProof}{\relax}
\providecommand{\FalseCl}{\bot}
\providecommand{\red}{\mathrel{\rightarrow}}
\providecommand{\varAssign}{\sigma}
\providecommand{\defiff}{\mathrel{\leftrightarrow}}
\providecommand{\FalseInt}{\bot}
\providecommand{\fn}{\operatorname{fn}}
\providecommand{\dotted}{\cdot}

% --- tag-macro semantic fix (2026-10-13, final): \iftag CHOOSES a branch and the OLP proof body is
% always the SECOND group (\iftag{tag}{body}, {tag}{body}{c}, {tag}{body}{c}{d}); a definition that
% prints the tag literal drops the proof body. Bodies contain blank lines, hence \long.
\makeatletter
\providecommand{\iftag}{}
\long\def\iftag#1#2{\@ifnextchar\bgroup{\iftag@three{#2}}{\iftag@stop{#2}}}
\long\def\iftag@three#1#2{\iftag@stop{#1}\@ifnextchar\bgroup{\iftag@eat}{\relax}}
\long\def\iftag@eat#1{}
\long\def\iftag@stop#1{#1}
\long\def\oliflabeldef#1#2{\@ifnextchar\bgroup{\@gobble}{\iftag@stop{#2}}}
\providecommand{\tagrefs}[1]{}
\providecommand{\settag}[1]{}
\makeatother
\begin{document}

% ==== OLP file: content/proof-theory/sequent-calculus/admissible-derivable.tex ====
% Part: proof-theory
% Chapter: sequent-calculus
% Section: admissible-derivable



\olfileid{pt}{seq}{adm}

\olsection{Admissible and Derivable Rules}

Many results of proof theory concern---or use---results about whether
everything that can be \olup{prove}d using a specific rule can also be
\olup{prove}d without that rule. The cut-elimination theorem is the most
well-known example. It shows that any sequent that can be \olup{prove}d
using the rules of a sequent system with~\Cut{} can also be \olup{prove}d
without~\Cut. Because this kind of property is so important, it has a
name.

\begin{defn}
A rule~$R$ is \emph{admissible} in a sequent system if any sequent
that can be proved in that system together with the rule~$R$ can also
be proved without~$R$.
\end{defn}

\begin{ex}\ollabel{ex:land-deriv}
The rule $\LeftR{\land}$ takes different forms in \Log{G1c} and
\Log{G3c}. In \Log{G1c} there are two versions of the rule, one with
the left conjunct $A$ of the principal \olup{formula}~$A \land B$ in
the premise, and one with the right conjunct~$B$:
\[
\Axiom$ A, \Gamma \fCenter \Delta$
\RightLabel{$\LeftR{\land}_1$}
\UnaryInf$ A \land B, \Gamma \fCenter \Delta$
\DisplayProof
\qquad
\Axiom$B, \Gamma \fCenter \Delta$
\RightLabel{$\LeftR{\land}_2$}
\UnaryInf$A \land B, \Gamma \fCenter \Delta$
\DisplayProof
\]
By contrast, \Log{G3c} has only one rule:
\[
\Axiom$ A, B, \Gamma \fCenter \Delta$
\RightLabel{$\LeftR{\land}_3$}
\UnaryInf$ A \land B, \Gamma \fCenter \Delta$
\DisplayProof
\]
We might now ask, ``Is \Log{G3c}'s rule $\LeftR{\land}_3$ admissible
in \Log{G1c}?'' It is: we can directly simulate any inference using
$\LeftR{\land}_3$ using only the rules of~\Log{G1c}. In addition to
$\LeftR{\land}_1$ and $\LeftR{\land}_2$ we'll also need the structural
rule~$\LeftR{\Contraction}$:
\[
\Axiom$A, B, \Gamma \fCenter \Delta$
\RightLabel{$\LeftR{\land}_1$}
\UnaryInf$A \land B, B, \Gamma \fCenter \Delta$
\RightLabel{$\LeftR{\land}_2$}
\UnaryInf$A \land B, A \land B, \Gamma \fCenter \Delta$
\RightLabel{\LeftR{\Contraction}}
\UnaryInf$A \land B \Gamma \fCenter \Delta$
\DisplayProof
\]
\end{ex}

Simulating a inferences using one rule using other rules is one way to
show that a rule is admissible. But it is not the only way, and some
rules are admissible but can't be so simulated. Rules that can be are
called \emph{derivable}.

\begin{defn}
A rule~$R$,
\[
\AxiomC{$S_1$}
\AxiomC{$\dots$}
\AxiomC{$S_n$}
\RightLabel{$R$}
\TrinaryInfC{$S$}
\DisplayProof
\]
is called \emph{derivable} in a system if there is a schematic
\olup{proof} of the sequent~$S$ using the rules and axioms of the system
together with additional initial sequents of the forms of~$S_1$, \dots,~$S_n$.
\end{defn}

The \olup{proof} given in \olref{ex:land-deriv} shows that
$\LeftR{\land}_3$ is a derivable rule in~\Log{G1c}, as that \olup{proof}
is a schematic \olup{proof} of the conclusion of $\LeftR{\land}_3$ from
the premises of $\LeftR{\land}_3$ as initial sequents using only rules
of \Log{G1c}. (It didn't use any of the axioms of \Log{G1c}.)

\begin{prop}\ollabel{prop:lor-G3-adm}
The rules $\LeftR{\land}$ and $\RightR{\lor}$ of~\Log{G3c} are
derivable in~\Log{G1c}.
\end{prop}

\begin{proof}
  The proof for $\LeftR{\land}$ is given in \olref{ex:land-deriv}. The
  proof for $\RightR{\lor}$ is left as an exercise.
\end{proof}

\begin{prob}
Show that the rule $\RightR{\lor}$ of \Log{G3c} is derivable
in~\Log{G1c}.
\end{prob}

\begin{prob}
Suppose $\liff$ is a defined \olup{operator}: $A \liff B$ is short for
$(A \lif B) \land (B \lif A)$. Show that the rules
\begin{enumerate}
\item \Axiom$\Gamma, A, B \fCenter \Delta$
\Axiom$\Gamma \fCenter \Delta, A, B$
\RightLabel{\LeftR{\liff}}
\BinaryInf$\Gamma \fCenter \Delta, A \liff B$
\DisplayProof
\item
\Axiom$A, \Gamma \fCenter \Delta, B$
\Axiom$B, \Gamma \fCenter \Delta, A$
\RightLabel{\RightR{\liff}}
\BinaryInf$\Gamma \fCenter \Delta, A \liff B$
\DisplayProof
\end{enumerate}
are derivable in (a)~\Log{G1c} and (b)~\Log{G3c}.
\end{prob}


\begin{ex}\ollabel{ex:weak-adm}
The rule $\RightR{\Weakening}$ of \Log{G1c},
\[
\Axiom$ \Gamma \fCenter \Delta$
\RightLabel{\RightR{\Weakening}}
\UnaryInf$ \Gamma \fCenter \Delta, A$
\DisplayProof
\]
is admissible
in~\Log{G3c}. The idea of the proof is simple: suppose there is
\olup{proof}~$\pi$ of the premise $\Gamma \Sequent \Delta$ in~\Log{G3c}.
Add the \olup{formula}~$A$ to the succedent of every sequent
in~$\pi$. Axioms remain axioms; correct inferences are still correct
inferences. The end-\olup{formula} of the new \olup{proof} is $\Gamma \Sequent
\Delta, A$.

The rule $\RightR{\Weakening}$ is, however, not derivable
in~\Log{G3c}. For suppose it were, i.e., suppose one could find
\olup{proof} of $\Gamma \Sequent
\Delta, A$ using only the axioms and rules of~\Log{G3c} plus possibly
the additional initial sequent~$\Gamma \Sequent \Delta$. If this
\olup{proof} is, as would be required for $\RightR{\Weakening}$ to be
derivable, purely schematic, we could turn it into \olup{proof}
of~$\Sequent A$ by deleting $\Gamma$ and~$\Delta$. This would turn
the additional initial sequents~$\Gamma \Sequent \Delta$ into the
empty sequent~$\quad \Sequent \quad$. Since the empty sequent contains
no \olup{formula}s at all, but all rules of~\Log{G3c} require at least
one active \olup{formula} in their premises, the empty sequent can't be a
premise of any inference in this schematic \olup{proof}. Consequently, the
\olup{proof} can only be schematic if it doesn't actually use the additional
sequent $\Gamma \Sequent \Delta$ as an initial sequent. In other
words, it would be a schematic derivation of $\Gamma \Sequent
\Delta, A$ using only the axioms and rules of~\Log{G3c} and we would
have shown that any sequent of this form has \olup{proof} in~\Log{G3c}.
But this is impossible, since \Log{G3c} is sound and so can only \olup{prove}
sequents that are true in every \olup{structure}. Taking $\Gamma = \Delta
= \emptyset$ and $A \ident \lfalse$, \Log{G3c} would have to \olup{prove},
e.g., the sequent $\quad \Sequent \lfalse$, which it does not.
\end{ex}

Let's give a nice inductive proof of the result that weakening is an
admissible rule in~\Log{G3c}. In fact, a slightly stronger result
holds, as the intuitive argument shows: by adding $A$ to the right side
of every sequent of~$\pi$, we do not change the structure of the
\olup{proof}, in particular, we do not increase the height. Since this is
important, it deserves its own definition.

\begin{defn}
A rule~$R$,
\[
\AxiomC{$S_1$}
\AxiomC{$\dots$}
\AxiomC{$S_n$}
\RightLabel{$R$}
\TrinaryInfC{$S$}
\DisplayProof
\]
is \emph{\olup{height}-preserving admissible} (or
\emph{\olup{hp}-admissible}) in a system if, whenever $\Proves[h_i] S_i$
for $i = 1$, \dots,~$n$, then $\Proves_m S$ with $m = \max(h_1, \dots,
h_n)$.
\end{defn}

Note that it is enough for a rule to be \olup{height}-preserving
admissable that, whenever the premises~$S_i$ have \olup{proof}s~$\pi_i$
with \olup{height}~$h_i = \pheight{\pi_i}$, there is \olup{proof}~$\pi$ of
the conclusion of \olup{height}~$\pheight{\pi}$ no greater than the
maximum of the~$h_i$. In particular, if there is only one
premise~$S_1$, it is enough that $\pheight{\pi} \le \pheight{\pi_1}$.
In the case of $\RightR{\Weakening}$ and $\LeftR{\Weakening}$ we
actually have $\pheight{\pi} = \pheight{\pi_1}$, but that is not
necessary.

\begin{prop}\ollabel{prop:weak-G3c-adm}
The rules $\RightR{\Weakening}$ and $\LeftR{\Weakening}$ of \Log{G1c} are
\olup{height}-preserving admissible in~\Log{G3c}.
\end{prop}

\begin{proof}
  We have to show that if $\Log{G3c} \Proves \Gamma \Sequent \Delta$
  with a \olup{proof}~$\pi$ of \olup{height}~$\pheight{\pi} = n$ then there
  is a \Log{G3c}-\olup{proof} $\pi'$ of $\Gamma \Sequent \Delta, A$ with
  $\pheight{\pi'} = n$. We prove this by induction on~$n$. If $n=0$,
  then $\pi$ is just the axiom $\Gamma \Sequent \Delta$ (i.e., some
  atomic $C$ occurs in both $\Gamma$ and~$\Delta$), and so is $\Gamma
  \Sequent \Delta, A$. If $n = \pheight{\pi} >0$, the \olup{proof}~$\pi$
  has a last inference. We distinguish cases according to which rule
  was used. We give just one example: suppose the last inference was
  $\RightR{\land}$. Then $\Delta = \Delta', B \land C$ and the
  immediate sub-\olup{proof}s $\pi_1$ and~$\pi_2$ of the last inference
  have end-sequents $\Gamma \Sequent \Delta', B$ and $\Gamma \Sequent
  \Delta', C$, respectively. In other words, $\pi$ looks like this:
  \[
  \AxiomC{}
  \RightLabel{$\pi_1$}
  \Deduce$\Gamma \fCenter \Delta', B$
  \AxiomC{}
  \RightLabel{$\pi_2$}
  \Deduce$\Gamma \fCenter \Delta', C$
  \RightLabel{\RightR{\land}}
  \BinaryInf$\Gamma \fCenter \Delta', B \land C$
  \DisplayProof
  \]
  We also have that $n =
  \max(\pheight{\pi_1}, \pheight{\pi_2}) + 1$. So $\pheight{\pi_1} <
  n$ and $\pheight{\pi_2} < n$, and the inductive hypothesis applies:
  there are \olup{proof}s $\pi_1'$ and~$\pi_2'$ of $\Gamma \Sequent
  \Delta', B, A$ and $\Gamma \Sequent \Delta', C, A$. Using the
  $\RightR{\land}$ rule, we obtain \olup{proof}~$\pi'$ of $\Gamma
  \Sequent \Delta', B \land C, A$:
  \[
  \AxiomC{}
  \RightLabel{$\pi_1'$}
  \Deduce$\Gamma \fCenter \Delta', B, A$
  \AxiomC{}
  \RightLabel{$\pi_2'$}
  \Deduce$\Gamma \fCenter \Delta', C, A$
  \RightLabel{\RightR{\land}}
  \BinaryInf$\Gamma \fCenter \Delta', B \land C, A$
  \DisplayProof
  \]
  We have $\pheight{\pi'} =
  \max(\pheight{\pi_1'}, \pheight{\pi_2'}) + 1 = \max(\pheight{\pi_1},
  \pheight{\pi_2}) + 1 = \pheight{\pi}$.
\end{proof}

\olref{prop:weak-G3c-adm} is extremely useful. We'll introduce a
notation for repeated application of it: if $\pi$ is \olup{proof}
of~$\Gamma \Sequent \Delta$ then $\pi[\Pi \Sequent \Lambda]$ is the
result of applying admissibility of weakening on the left with all the
\olup{formula}s in~$\Pi$ and admissibility of weakening on the right with
all \olup{formula}s in~$\Lambda$. It is \olup{proof} of $\Gamma, \Pi
\Sequent \Delta, \Lambda$.

% ==== OLP file: content/proof-theory/sequent-calculus/interpretation-rules.tex ====
% Part: proof-theory
% Chapter: sequent-calculus
% Section: interpretation-rules



\olfileid{pt}{seq}{irl}

\olsection{Interpretation of Rules}

Recall the interpretation of a sequent: A sequent $\Gamma \Sequent
\Delta$ is true in \olup{structure}~$\Struct{M}$ if either at least one
$A \in \Gamma$ is false or at least one~$A \in \Delta$ is true. We
can read the logical rules of~\Log{G3c} as codifying the truth
conditions for their principal \olup{formula}s. Consider \RightR{\lif}.
Its principal \olup{formula} is $A \lif B$. It is true if either $A$
is false or~$B$ is true. Hence the sequent $\ \Sequent A \lif B$ is
true iff $A \Sequent B$ is true. If we add the context \olup{formula}s
$\Gamma$, $\Delta$ to the antecedents and succedents of these
sequents, we get exactly the conclusion and premise of the
\RightR{\lif} rule. The situation is similar for~\LeftR{\lif}. $A
\lif B$ is false if $A$~is true and $B$~is false. So $A \lif B
\Sequent \ $ is true iff both $\ \Sequent A$ and $B \Sequent \ $ are
true. The conclusion of \LeftR{\lif} is equivalent to the
conjunction of its premises. This pattern ensures that the rules
of~\Log{G3c} are sound (if all premises are true, the conclusion is
true). They also ensure completeness.

In fact, the truth conditions for any truth-functional connective can
be described in this way as a set of sequents. This follows from the
fact that every truth-function in classical logic can be described by
\olup{formula} in conjunctive normal form: a conjunction of disjunctions
of atomic and negated atomic \olup{formula}s. For instance, take $A
\oplus B$ to be true iff exactly one of $A$ and $B$ is true, i.e.,
``exclusive or.'' Then $A \oplus B$ is true iff $(A \lor B) \land
(\lnot A \lor \lnot B)$ is, and it is false if $(\lnot A \lor B)
\land (A \lor \lnot B)$ is true. Hence, we could introduce sequent
calculus rules for $\oplus$ as
\begin{align*}
&
\Axiom$A, \Gamma \fCenter \Delta, B$
\Axiom$B, \Gamma \fCenter \Delta, A$
\RightLabel{\LeftR{\oplus}}
\BinaryInf$A \oplus B, \Gamma \fCenter \Delta$
\DisplayProof
\\
&
\Axiom$\Gamma \fCenter \Delta, A, B$
\Axiom$A, B, \Gamma \fCenter \Delta$
\RightLabel{\RightR{\oplus}}
\BinaryInf$\Gamma \fCenter \Delta, A \oplus B$
\DisplayProof
\end{align*}
and these rules would also be sound and complete. Their conclusions
are true iff all premises are true.

\begin{prob}
Suppose $A \downarrow B$ is true iff neither $A$ nor~$B$ is true
(``NOR''). What would its \Log{G3c} rules be? (Find two sequents which
are true at the same time iff $A \downarrow B$ is true. That gives
you the \RightR{\downarrow} rule. Then find one sequent that is true
iff $A \downarrow B$ is false. That gives you the \LeftR{\downarrow}
rule.)
\end{prob}

In \Log{G3c}, each connective and quantifier has exactly two rules: a
left rule and a right rule. In \Log{G1c}, however, there are two
\LeftR{\land} and two \RightR{\lor} rules. The reason is that
Gentzen's original system~\Log{LK}, on which \Log{G1c} is modeled, had
a variant~\Log{LJ} for an important non-classical logic, viz.,
intuitionistic logic. To get a system that was sound and complete for
intuitionistic logic, the succedents of any sequent in~\Log{LJ} is
restricted to contain at most one \olup{formula}. This makes it
impossible to use the \RightR{\lor} rule of~\Log{G3c}, as its premise
contains both $A$ and~$B$ in the succedent. Gentzen thus used a pair
of \RightR{\lor} rules for \Log{LJ}, and used the same for~\Log{LK}.
The intuitionistic version~\Log{G1i} similarly is just~\Log{G1c} with
all succedents restricted to at most one \olup{formula}. (\Log{G3c} also
has an intuitionistic version, but it some of its rules differ from
the corresponding rules in~\Log{G3c}. For instance, it, too, uses a
pair of \RightR{\lor} rules.)

That the structural rules of~\Log{G1c} are sound is easy to see. For
instance, if $\Gamma \Sequent \Delta$ contains a false \olup{formula} on
the right or a true \olup{formula} on the left, so does $\Gamma \Sequent
\Delta, A$: the same one, whether $A$ is true or false is
irrelevant. So if the premise of~\RightR{\Weakening} is true, so is
its conclusion. Note that here the converse is not true. The
conclusion may be true because $A$ is true, and then the premise is
not guaranteed to be true.

Why have structural rules at all? Contraction (\RightR{\Contraction},
\LeftR{\Contraction}) is needed for instance to \olup{prove} $\ \Sequent A
\lor \lnot A$. If you start with $A \Sequent A$ you first get
$\Sequent A, \lnot A$ by~\RightR{\lnot}. Then you can use the two
versions of \RightR{\lor} in succession, once to get $A \lor \lnot
A$ from $\lnot A$, and once to get $A \lor \lnot A$ from~$A$.
What you get is $\ \Sequent A \lor \lnot A, A \lor \lnot A$. So to
get $\Sequent A \lor \lnot A$ in~\Log{G1c} you need to be able to
``contract'' two occurrences of the same \olup{formula} in \olup{proof}:
\[
\Axiom$A \fCenter A$
\RightLabel{\RightR{\lnot}}
\UnaryInf$\fCenter A, \lnot A$
\RightLabel{\RightR{\lor}}
\UnaryInf$\fCenter A, A \lor \lnot A$
\RightLabel{\RightR{\lor}}
\UnaryInf$\fCenter A \lor \lnot A, A \lor \lnot A$
\RightLabel{\RightR{\Contraction}}
\UnaryInf$\fCenter A \lor \lnot A$
\DisplayProof
\]
There are in fact valid sequents which you could not \olup{prove}
in~\Log{G1c} without weakening or contraction. For instance,
\olup{proof} of the sequent $A \Sequent B \lif A$ in~\Log{G1c} has to
end in \RightR{\lif} (since only $B \lif A$ can be principal), and
that requires $B, A \Sequent A$ as a premise. This, in turn,
requires \LeftR{\Weakening} to get it from the axiom $A \Sequent A$:
\[
\Axiom$A \fCenter A$
\RightLabel{\LeftR{\Weakening}}
\UnaryInf$B, A \fCenter A$
\RightLabel{\RightR{\lif}}
\UnaryInf$A \fCenter B \lif A$
\DisplayProof
\]

In \Log{G3c} the contraction and weakening rules are not needed at
all. Weakening is built into the axioms. Contraction can mostly be
avoided by merging the two copies of the context $\Gamma$, $\Delta$ in
rules with two premises into one. Both \Log{G1c} and \Log{G3c} have
such ``context sharing'' rules. In one-premise rules, contraction is
only needed for \RightR{\lor}, \LeftR{\land}, \RightR{\lexists},
and~\LeftR{\lforall}. The versions of these rules in~\Log{G3c} make
contraction entirely superfluous. There is a system \Log{G2c} (see
\olref{tab:G2c}) in which the two-premise rules have independent
contexts. In \Log{G2c} contraction is even more important than
in~\Log{G1c}.

\subfile{rules-G2c}

% ==== OLP file: content/proof-theory/sequent-calculus/introduction.tex ====
% Part: proof-theory
% Chapter: sequent-calculus
% Section: introduction



\olfileid{pt}{seq}{int}

\olsection{Introduction}

Sequent calculi are \olup{proof} systems that were introduced in the
1930s by Gerhard Gentzen. They were intended less as a practical
method for \olup{proving} things. Rather, Gentzen used them in order to
be able to prove certain results about the possible structure of
\olup{proof}s and their property. These kinds of questions are
exactly the questions with which proof theory is concerned.

Sequent calculi operate on \emph{sequents}, as the name suggests, and
not on \olup{formula}s. A sequent is a pair of sequences or multisets of
\olup{formula}s. Gentzen's original systems (\Log{LK} and \Log{LJ}) used
sequences; proof theorists now prefer to use multisets. A multiset is
like a set, except that a multiset can contain an \olup{element} more
than once. Just as in sets, the order of \olup{element}s does not matter,
however. So the expressions $\{A, A, B\}$ and $\{A, B, A\}$
describe the same multiset, but that multiset is different from $\{A,
B\}$ and from $\{A, A, B, B\}$, since the latter contain
different numbers of $A$ and $B$ than the former.

Traditionally, multisets of \olup{formula}s in proof theory are written
using uppercase Greek letters. So when we write $\Gamma$, $\Delta$,
$\Pi$, $\Lambda$, or $\Theta$, we intend these to be multisets of
\olup{formula}s of whatever logic we are working with. Traditionally
also, instead of writing $\tuple{\Gamma, \Delta}$ proof theorists
separate the two components by a sequent symbol; we use~$\Sequent$. 

\begin{defn}[Sequent]
A \emph{sequent} is an expression of the form
\[
\Gamma \Sequent \Delta
\]
where $\Gamma$ and $\Delta$ are finite (possibly empty) multisets of
\olup{formula}s. $\Gamma$ is called the \emph{antecedent}, while
$\Delta$ is the \emph{succedent}.
\end{defn}

Let $G$~be the conjunction of the \olup{formula}s in~$\Gamma$ (taking
the empty conjunction to be~$\ltrue$) and $D$~be the disjunction of
the \olup{formula}s in~$\Delta$ (taking the empty disjunction to
be~$\lfalse$). Then the sequent $\Gamma \Sequent \Delta$ is true (in
some structure~$\Struct{M}$) iff $G \lif D$ is true. In classical
logic, this amounts to the same as saying that either one of the
\olup{formula}s in~$\Gamma$ is false, or one of the \olup{formula}s
in~$\Delta$ is true.

If $\Gamma$ is a multiset of \olup{formula}s, we write $\Gamma, A$ (or
$A, \Gamma$) for the result of adding $A$ to~$\Gamma$. If $\Delta$
is also a multiset of \olup{formula}s, then $\Gamma, \Delta$ is the
multiset union of the two sequences---if $\Gamma$ contains $A$ with
\emph{multiplicity}~n (i.e., it contains $n$ copies of~$A$) and
$\Delta$ with multiplicity~$m$, then $\Gamma, \Delta$ contains $A$
with multiplicity~$n+m$. If the multiset on one side of a sequent is
empty, we usually just leave a blank. For instance, $\Sequent A$ is
the sequent with an empty antecendent and the succedent contains $A$
with multiplicity~$1$.

!!^a{proof} in the sequent calculus is a tree of sequents, where
the topmost (or initial) sequents are axioms in the system considered,
and if a sequent stands below one or two other sequents, it must
follow correctly by a rule of inference in the system. We will first
consider two different sequent systems for classical logic,
$\Log{G1c}$ and~$\Log{G3c}$. Their axioms and rules are given in
\cref{pt:seq:int:tab:G1c,pt:seq:int:tab:G3c}. \oliflabeldef{fol:seq::chap}{ (Gentzen's
original system \Log{LK} is described in \olref[fol][seq][]{chap}.)}{}

\subfile{rules-G1c}
\subfile{rules-G3c}

The systems are subtly different, and these differences mean that they
have different proof theoretic properties. \Log{G1c} has axioms $A
\Sequent A$ with no other \olup{formula}s allowed. \Log{G3c} has axioms
of the form $C, \Gamma \Sequent \Delta, C$ with arbitrary ``side''
\olup{formula}s in $\Gamma$ and~$\Delta$, but the \olup{formula}~$C$
appearing on both sides must be \emph{atomic}. \Log{G1c} has two versions
each of the $\LeftR{\land}$ and $\RightR{\lor}$ rules whereas \Log{G3c}
has just one each. And \Log{G1c} has additional ``structural rules''
called ``contraction'' (\LeftR{\Contraction}, \RightR{\Contraction})
and ``weakening'' (\LeftR{\Weakening}, \RightR{\Weakening}).

The systems \Log{G1c} and \Log{G3c} happen to be sound and complete
for classical logic. In other words, every sequent \olup{provable} in
these systems is true in every \olup{structure}, and every sequent that
is true in every \olup{structure} has \olup{proof}. So the question
\emph{whether} a sequent has \olup{proof} can also be answered by other
methods of logic (e.g., those of model theory). And since both systems
\olup{prove} exactly the sequents true in every \olup{structure}, they in
particular prove the same sequents.

Proof theory, however, is not just concerned with \emph{whether}
something is provable, but \emph{how}. Proof theorists consider
questions such as ``Can a certain rule always be avoided?,'' ``Can
this rule be added to the system without resulting in new sequents
becoming provable?,'' or ``Can \olup{proof}s in one system be turned into
\olup{proof}s in another?'' If the answer is ``yes'' to such a question,
the proof of this fact is often proved by induction on the complexity
of \olup{proof}s, or, equivalently, by induction on the structure of
\olup{proof}s. This yields not just a proof of the fact (e.g., if
\Log{G1c} proves a sequent then \Log{G3c} proves the same sequent),
but a procedure (e.g., of turning \olup{proof}s in~\Log{G1c} into
\olup{proof}s in~\Log{G3c}).

A central result in proof theory is the \emph{cut-elimination
theorem}. It shows that we can add the cut rule
\[
\Axiom$\Gamma \fCenter \Delta, A$
\Axiom$A, \Pi \fCenter \Lambda$
\RightLabel{\Cut}
\BinaryInf$\Gamma, \Pi \fCenter \Delta, \Lambda$
\DisplayProof
\]
to the sequent systems without changing which sequents are
\olup{provable}. Moreover, the proof provides a method that transforms
any \olup{proof} using the rules of \Log{G1c} (or \Log{G3c}) together with
\Cut{} into one that only uses the rules of \Log{G1c} (or \Log{G3c}).
In other words, this procedure ``eliminates'' the \Cut{} rule from
\olup{proof}s using it; hence the name.


% ==== OLP file: content/proof-theory/sequent-calculus/invertibility.tex ====
% Part: proof-theory
% Chapter: sequent-calculus
% Section: invertibility



\olfileid{pt}{seq}{inv}

\olsection{Invertibility of Rules}

\begin{defn}
A rule~$R$,
\[
\AxiomC{$S_1$}
\AxiomC{$\dots$}
\AxiomC{$S_n$}
\RightLabel{$R$}
\TrinaryInfC{$S$}
\DisplayProof
\]
is \emph{invertible} in a system if, whenever $\Proves S$ then
$\Proves[h_i] S_i$ for $i = 1$, \dots,~$n$. It is
\emph{\olup{height}-preserving invertible} (or \emph{\olup{hp}-invertible})
in a system if, whenever $\Proves_n S$ then $\Proves[n] S_i$ for $i =
1$, \dots,~$n$.
\end{defn}

\begin{lem}\ollabel{lem:G3c-invert}
The rules of~\Log{G3c} for $\lnot$, $\land$, $\lor$, and $\lif$ are
\olup{height}-preserving invertible, that is:
\begin{enumerate}
  \item \RightR{\lnot}: If $\Log{G3c} \Proves[n] \Gamma \Sequent
  \Delta, \lnot A$, then $\Log{G3c} \Proves[n] A, \Gamma \Sequent
  \Delta$.
  \item \LeftR{\lnot}: If $\Log{G3c} \Proves[n] \Gamma, \lnot A
  \Sequent \Delta$, then $\Log{G3c} \Proves[n] \Gamma \Sequent \Delta,
  A$.
  \item \RightR{\land}: If $\Log{G3c} \Proves[n] \Gamma \Sequent
  \Delta, A \land B$, then $\Log{G3c} \Proves[n] \Gamma \Sequent
  \Delta, A$ and $\Log{G3c} \Proves[n] \Gamma \Sequent
  \Delta, B$.
  \item \LeftR{\land}: If $\Log{G3c} \Proves[n] \Gamma, A \land B
  \Sequent \Delta$, then $\Log{G3c} \Proves[n] A, B, \Gamma \Sequent
  \Delta$.
  \item \RightR{\lor}: If $\Log{G3c} \Proves[n] \Gamma \Sequent
  \Delta, A \lor B$, then $\Log{G3c} \Proves[n] \Gamma \Sequent
  \Delta, A, B$.
  \item \LeftR{\lor}: If $\Log{G3c} \Proves[n] A \lor B, \Gamma \Sequent
  \Delta$, then $\Log{G3c} \Proves[n] A, \Gamma \Sequent
  \Delta$ and $\Log{G3c} \Proves[n] B, \Gamma \Sequent
  \Delta$.
  \item \RightR{\lif}: If $\Log{G3c} \Proves[n] \Gamma \Sequent
  \Delta, A \lif B$, then $\Log{G3c} \Proves[n] A, \Gamma \Sequent
  \Delta, B$.
  \item \LeftR{\lif}: If $\Log{G3c} \Proves[n] A \lif B, \Gamma
  \Sequent \Delta$, then $\Log{G3c} \Proves[n] \Gamma \Sequent \Delta,
  A$ and $\Log{G3c} \Proves[n] B, \Gamma \Sequent \Delta$.
\end{enumerate}
\end{lem}

\begin{proof}
We have to show that, for any rule~$R$ of~\Log{G3c}, if there is
\olup{proof}~$\pi$ of the conclusion~$S$ of~$R$, then there are also
\olup{proof}s~$\pi_i$ of the premises $S_i$ of~$R$, with $\pheight{\pi_i}
\le \pheight{\pi}$. Consider the rule \LeftR{\land}: suppose there is
\olup{proof}~$\pi$ with end-sequent of the form $A \land B, \Gamma
\Sequent \Delta$. We have to show that there is \olup{proof} $\pi'$ of
$A, B, \Gamma \Sequent \Delta$ with $\pheight{\pi'} \le
\pheight{\pi}$. We do this by induction on $n = \pheight{\pi}$.

If $n = 0$, $\pi$ consists of just an axiom: $A \land B, \Gamma
\Sequent \Delta$. In \Log{G3c}, axioms are of the form $C, \Gamma
\Sequent \Delta, C$ with $C$ atomic. In other words, $C$ can't be
$A \land B$. If $A \land B, \Gamma \Sequent \Delta$ is an axiom,
then some atomic $C$ must be \olup{element} of both $\Gamma$
and~$\Delta$. But then $A, B, \Gamma \Sequent \Delta$ is also an
axiom, and we have found our \olup{proof}~$\pi'$ (which is also of
\olup{height}~$0$).

If $n>0$ it ends in an inference. We have to prove the claim for~$n$,
but may assume that it is true for any \olup{proof} of height~$<n$, even
if the context is different. In other words, for any $k<n$ and any
$\Pi$ and $\Lambda$, if $A \land B, \Pi \Sequent \Lambda$ has
\olup{proof} of \olup{height}~$k$, then there is \olup{proof} of
\olup{height}~$\le k$ of $A, B, \Pi \Sequent \Lambda$.

We have two cases: either $A \land
B$ on the left is principal, or it is not. If it is principal, the
last inference is \LeftR{\land} and the immediate
sub-\olup{proof}~$\pi_1$ ends in $A, B, \Gamma \Sequent \Delta$. The
result holds in this case as $\pheight{\pi_1} < \pheight{\pi}$.

If $A \land B$ is not principal, some other \olup{formula} is and $A
\land B$ belongs to the context of the last inference. Then the
inductive hypothesis applies to the immediate sub-\olup{proof}s of the
last rule~$R$, which yields \olup{proof}~$\pi_1'$ or two
\olup{proof}s~$\pi_1'$ and~$\pi_2'$ of \olup{height} no greater than the
height(s) of the immediate sub-\olup{proof}s of~$\pi$. The end-sequents
of $\pi_1'$ and $\pi_2'$ contain $A, B$ in the context on the left
instead of $A \land B$. Apply the rule~$R$ to obtain the
proof~$\pi'$. For instance, suppose $\pi$ ends in~\LeftR{\lif}:
\[
\AxiomC{}
\RightLabel{$\pi_1$}
\Deduce$A \land B, \Gamma' \fCenter \Delta, C$
\AxiomC{}
\RightLabel{$\pi_2$}
\Deduce$A \land B, \Gamma', D \fCenter \Delta$
\RightLabel{\LeftR{\lif}}
\BinaryInf$A \land B, \Gamma', C \lif D \fCenter \Delta$
\DisplayProof
\]
Here $C \lif D$ on the left is principal, and the context on the
left is $A \land B, \Gamma'$ (i.e., $\Gamma = \Gamma', C \lif D$).
The inductive hypothesis yields \olup{proof}s $\pi_1'$ and $\pi_2'$ of 
$A \land B, \Gamma' \Sequent \Delta, C$ and $A \land B, \Gamma',
D \Sequent \Delta$, respectively. These two sequents can again
serve as premises for the \LeftR{\lif} rule to get~$\pi'$:
\[
\AxiomC{}
\RightLabel{$\pi_1'$}
\Deduce$A, B, \Gamma' \fCenter \Delta, C$
\AxiomC{}
\RightLabel{$\pi_2'$}
\Deduce$A, B, \Gamma', D \fCenter \Delta$
\RightLabel{\LeftR{\lif}}
\BinaryInf$A, B, \Gamma', C \lif D \fCenter \Delta$
\DisplayProof
\]
We have 
\begin{align*}
\pheight{\pi} & = \max(\pheight{\pi_1}, \pheight{\pi_2}) + 1\\
\pheight{\pi'} & = \max(\pheight{\pi_1'}, \pheight{\pi_2'}) + 1
\intertext{by definition, and}
\pheight{\pi_1'} & \le \pheight{\pi_1}\\
\pheight{\pi_2'} & \le \pheight{\pi_2}
\intertext{by inductive hypothesis. So,}
\pheight{\pi'} & \le \pheight{\pi}.
\end{align*}
\end{proof}

\begin{prob}
Give the proof of \olref{lem:G3c-invert} for \RightR{\lif}.
\end{prob}

\begin{lem}\ollabel{lem:invert-quant}
The rules \RightR{\lforall} and \LeftR{\lexists} are
\olup{height}-preserving invertible in the following sense:
\begin{enumerate}
  \item If $\Log{G3c} \Proves[n] \Gamma \Sequent \Delta,
  \lforall[x][A(x)]$ then $\Log{G3c} \Proves[n] \Gamma \Sequent
  \Delta, A(c)$, and
  \item if $\Log{G3c} \Proves[n] \lexists[x][A(x)], \Gamma \Sequent
  \Delta$ then $\Log{G3c} \Proves[n] A(c), \Gamma \Sequent \Delta$,
\end{enumerate}
for any $c$ not occuring in $\Gamma$, $\Delta$, and $\lforall[x][A(x)]$
or $\lexists[x][A(x)]$, respectively.
\end{lem}

\begin{proof}
  We show (1) and leave (2) as an exercise. The argument is similar to
  the proof of \olref{lem:G3c-invert}. In the inductive step, there
  are again two cases. The case where $\lforall[x][A(x)]$ is
  principal again is quick: the premise of the last
  $\RightR{\lforall}$ inference is of the required form $\Gamma
  \Sequent \Delta, A(a)$ and the immediate sub-\olup{proof}~$\pi_1$
  ending in it has $\pheight{\pi_1} < \pheight{\pi}$. Moreover, the
  eigenvariable condition gives us that $a$ does not occur in
  $\Gamma$, $\Delta$, or $\lforall[x][A(x)]$. Since we may assume that
  $\pi_1$ is regular, $\Subst{\pi_1}{c}{a}$ is \olup{proof} of the same
  \olup{height} of $\Gamma \Sequent \Delta, A(c)$ by
  \olref[qua]{lem:subst-regular}.

  In the second case, $\lforall[x][A(x)]$ is not principal in the
  last inference; some other formula is. We again consider a case by
  way of example. Suppose the last inference is~\RightR{\land}. Some
  \olup{formula} in $\Delta$ is of the form~$B \land C$ and is principal.
  The \olup{proof}~$\pi$ ends in
  \[
\AxiomC{}
\RightLabel{$\pi_1$}
\Deduce$\Gamma \fCenter \Delta', \lforall[x][A(x)], B$
\AxiomC{}
\RightLabel{$\pi_2$}
\Deduce$\Gamma \fCenter \Delta', \lforall[x][A(x)], C$
\RightLabel{\RightR{\land}}
\BinaryInf$\Gamma \fCenter \Delta', \lforall[x][A(x)], B \land C$
\DisplayProof
\]
where $\Delta = \Delta', B \land C$. Let $c$ be \olup{constant} that
does not appear in $\Delta$; then clearly it also does not occur in
$\Delta', B$ or in $\Delta', C$. So the inductive hypothesis applies
to $\pi_1$ and~$\pi_2$; we have \olup{proof}s~$\pi_1'$ and $\pi_2'$ with
$\pheight{\pi_1'} \le \pheight{\pi_1}$ and $\pheight{\pi_2'} \le
\pheight{\pi_2}$. The required \olup{proof}~$\pi'$ of $\Gamma \Sequent
\Delta, A(c)$ is:
\[
\AxiomC{}
\RightLabel{$\pi_1'$}
\Deduce$\Gamma \fCenter \Delta', A(c), B$
\AxiomC{}
\RightLabel{$\pi_2'$}
\Deduce$\Gamma \fCenter \Delta', A(c), C$
\RightLabel{\RightR{\land}}
\BinaryInf$\Gamma \fCenter \Delta', A(c), B \land C$
\DisplayProof
\]
with $\pheight{\pi'} = \max(\pheight{\pi_1'}, \pheight{\pi_2'})+1 \le \pheight{\pi}$.
\end{proof}

\Olref{lem:G3c-invert} plays an important role in the proof of the
cut-elimination theorem. But it can also be used to show the
following:

\begin{prop}\ollabel{prop:G3c-cont-adm}
The contraction rules \LeftR{\Contraction} and \RightR{\Contraction}
of~\Log{G1c} are \olup{height}-preserving admissible in~\Log{G3c}.
\end{prop}

\begin{proof}
We give the argument for \RightR{\Contraction} and
\LeftR{\Contraction} at the same time. Suppose $\pi$ is \olup{proof} of
$\Gamma \Sequent \Delta, A, A$ or of $A, A, \Gamma \Sequent
\Delta$ in~\Log{G3c}. We have to show that there is also a
\Log{G3c}-\olup{proof}~$\pi'$ of $\Gamma \Sequent \Delta, A$ or $A,
\Gamma \Sequent \Delta$, respectively, with $\pheight{\pi'} \le
\pheight{\pi}$. We do this by induction on $n = \pheight{\pi}$.

If $n=0$, $\pi$ is just an axiom. But if $\Gamma \Sequent \Delta, A,
A$ is an axiom of~\Log{G3c}, so is $\Gamma \Sequent \Delta, A$:
either some atomic~$C$ is \olup{element} of both $\Gamma$ and~$\Delta$,
or $A$ is atomic and \olup{element} of~$\Gamma$. The same reasoning
applies if the end-sequent is $A, A, \Gamma \Sequent
\Delta$.

If $n>0$, $\pi$ ends in some rule~$R$. If $A$ is not principal in
this last inference, we can obtain \olup{proof}~$\pi'$ as in the proof
of \olref{lem:G3c-invert} by applying the inductive hypothesis to the
immediate sub-\olup{proof}s of~$\pi$ and applying the same rule~$R$ to
the resulting \olup{proof}(s). The inductive hypothesis is that if there
is \olup{proof} of \olup{height}~$k<n$ of either $\Pi \Sequent \Lambda, D,
D$ or $D, D, \Pi \Sequent \Lambda$, then there is \olup{proof} of
\olup{height}~$\le k$ of $\Pi \Sequent \Lambda, D$ or $D, \Pi \Sequent
\Lambda$, respectively. (Note that the inductive hypothesis applies if
the contexts or contracted \olup{formula} are different from $\Gamma$,
$\Delta$, or~$A$!)

For instance, suppose $R$ is $\RightR{\land}$
and $\pi$~ends in
\[
\AxiomC{}
\RightLabel{$\pi_1$}
\Deduce$\Gamma \fCenter \Delta', B, A, A$
\AxiomC{}
\RightLabel{$\pi_2$}
\Deduce$\Gamma \fCenter \Delta', C, A, A$
\RightLabel{\RightR{\land}}
\BinaryInf$\Gamma \fCenter \Delta', B \land C, A, A$
\DisplayProof
\]
where $\Delta = \Delta', B \land C$. The inductive hypothesis yields
\olup{proof}s $\pi_1'$ and $\pi_2'$ of $\Gamma \fCenter \Delta', B, A$
and $\Gamma \fCenter \Delta', C, A$, respectively, with
$\pheight{\pi_1'} \le  \pheight{\pi_1}$ and $\pheight{\pi_2'} \le
\pheight{\pi_2}$. The \olup{proof}~$\pi'$ is:
\[
\AxiomC{}
\RightLabel{$\pi_1'$}
\Deduce$\Gamma \fCenter \Delta', B, A$
\AxiomC{}
\RightLabel{$\pi_2'$}
\Deduce$\Gamma \fCenter \Delta', C, A$
\RightLabel{\RightR{\land}}
\BinaryInf$\Gamma \fCenter \Delta', B \land C, A$
\DisplayProof
\]

If $A$ is principal in the last inference, we apply the inversion
lemma (\olref{lem:G3c-invert}) to the immediate sub-\olup{proof}s
of~$\pi$, then the inductive hypothesis, and then rule~$R$ again. For
instance, suppose $\pi$ proves $\Gamma \Sequent \Delta, A, A$, $A
\ident (B \land C)$ and principal in the last inference. Then $\pi$
looks as follows:
\[
\AxiomC{}
\RightLabel{$\pi_1$}
\Deduce$\Gamma \fCenter \Delta, B \land C, B$
\AxiomC{}
\RightLabel{$\pi_2$}
\Deduce$\Gamma \fCenter \Delta, B \land C, C$
\RightLabel{\RightR{\land}}
\BinaryInf$\Gamma \fCenter \Delta, B \land C, B \land C$
\DisplayProof
\]
Apply \olref{lem:G3c-invert} to $\pi_1$ to get \olup{proof}~$\pi_1'$ of
$\Gamma \Sequent \Delta, B, B$. (It also yields \olup{proof} of
$\Gamma \Sequent \Delta, C, B$, but we don't need that.) Similarly,
by applying \olref{lem:G3c-invert} to~$\pi_2$ we obtain
\olup{proof}~$\pi_2'$ of $\Gamma \Sequent \Delta, C, C$. As the
inversion of \RightR{\land} is \olup{height}-preserving,
$\pheight{\pi_1'} \le  \pheight{\pi_1} < \pheight{\pi}$ and
$\pheight{\pi_2'} \le  \pheight{\pi_2} < \pheight{\pi}$. Thus, the
inductive hypothesis applies to $\pi_1'$ and $\pi_2'$: there are
\olup{proof}s~$\pi_1''$ and $\pi_2''$ of $\Gamma \Sequent \Delta, B$ and
$\Gamma \Sequent \Delta, C$, respectively, with $\pheight{\pi_1''}
\le \pheight{\pi_1'} \le \pheight{\pi_1}$ and $\pheight{\pi_2''} \le
\pheight{\pi_2'} \le \pheight{\pi_2}$. The \olup{proof}~$\pi'$ then is:
\[
\AxiomC{}
\RightLabel{$\pi_1''$}
\Deduce$\Gamma \fCenter \Delta, B$
\AxiomC{}
\RightLabel{$\pi_2''$}
\Deduce$\Gamma \fCenter \Delta, C$
\RightLabel{\RightR{\land}}
\BinaryInf$\Gamma \fCenter \Delta, B \land C$
\DisplayProof
\]
with $\pheight{\pi'} \le \pheight{\pi}$.

We must take special care for the quantifier rules in the cases where
the \olup{formula} is principal.

Suppose $A \ident \lforall[x][B(x)]$ is principal on the right. Then $\pi$ ends in
\[
\AxiomC{}
\RightLabel{$\pi_1$}
\Deduce$\Gamma \fCenter \Delta, \lforall[x][B(x)], B(a)$
\RightLabel{\RightR{\lexists}}
\UnaryInf$\Gamma \fCenter \Delta, \lforall[x][B(x)], \lforall[x][B(x)]$
\DisplayProof
\]
By \olref{lem:invert-quant}, there is \olup{proof}~$\pi_1'$ with
\olup{height} $\le \pheight{\pi_1}$ of $\Gamma \fCenter \Delta, B(c),
B(a)$ for any~$c$ not occuring in $\Gamma \fCenter \Delta,
\lforall[x][B(x)], B(a)$. We may assume that $\pi_1'$ is regular; so
by \olref[qua]{lem:subst-regular} the \olup{proof} $\Subst{\pi_1'}{a}{c}$
is \olup{proof} of $\Gamma \fCenter \Delta, B(a),
B(a)$ also of \olup{height} $\le \pheight{\pi_1}$. Now the inductive
hypothesis applies and gives us \olup{proof}~$\pi_1''$ of $\Gamma
\fCenter \Delta, B(a)$. Apply the \RightR{\lforall} rule to
get~$\pi'$:
\[
\AxiomC{}
\RightLabel{$\pi_1''$}
\Deduce$\Gamma \fCenter \Delta, B(a)$
\RightLabel{\RightR{\lforall}}
\UnaryInf$\Gamma \fCenter \Delta, \lforall[x][B(x)]$
\DisplayProof
\]
with $\pheight{\pi'} \le \pheight{\pi}$.

Now suppose $A \ident \lexists[x][B(x)]$
and principal on the right. Then $\pi$ looks as follows:
\[
\AxiomC{}
\RightLabel{$\pi_1$}
\Deduce$\Gamma \fCenter \Delta, \lexists[x][B(x)], \lexists[x][B(x)], B(t)$
\RightLabel{\RightR{\lexists}}
\UnaryInf$\Gamma \fCenter \Delta, \lexists[x][B(x)], \lexists[x][B(x)]$
\DisplayProof
\]
The inductive hypothesis applies to~$\pi_1$ and there is
\olup{proof}s~$\pi_1'$ of $\Gamma \Sequent \Delta, \lexists[x][B(x)],
B$. (Note that no inversion lemma is needed here!) The
\olup{proof}~$\pi'$ then is:
\[
\AxiomC{}
\RightLabel{$\pi_1'$}
\Deduce$\Gamma \fCenter \Delta,  \lexists[x][B(x)], B$
\RightLabel{\RightR{\lexists}}
\UnaryInf$\Gamma \fCenter \Delta, \lexists[x][B(x)]$
\DisplayProof
\]
with $\pheight{\pi'} \le \pheight{\pi}$.
\end{proof}

\begin{prob}
Sketch the proof of \olref{prop:G3c-cont-adm} for \LeftR{\Contraction}.
Describe the cases where $A \ident (B \land C)$ and $A \ident (B
\lif C)$.
\end{prob}

\begin{prob}
Sketch the proof of \olref{prop:G3c-cont-adm} for \LeftR{\Contraction}
for the remaining quantifier cases, i.e., where $A \ident
\lforall[x][B(x)]$ and $A \ident \lexists[x][B(x)]$.
\end{prob}


% ==== OLP file: content/proof-theory/sequent-calculus/proof-examples.tex ====
% Part: proof-theory
% Chapter: sequent-calculus
% Section: proof-examples



\olfileid{pt}{seq}{ex}

\olsection{Examples of Proofs}

To give \olup{proof}s of sequents it is typically convenient to work your
way upwards from the end-sequent by selecting \olup{formula} in it as
the active \olup{formula} and writing the premises of the corresponding
rule on top of the sequent, then repeating the process until it leads
to axioms. For instance, suppose we want to prove the sequent
\[(C \land D) \lif E \Sequent C \lif (D \lif E)\]
Consider $C \lif (D \lif E)$: it is of the form $A \lif B$ and
occurs on the right of the sequent. Matching the entire sequent with
the $\RightR{\lif}$ rule of \Log{G1c},
\[\Axiom$ A, \Gamma \fCenter \Delta, B$
\RightLabel{\RightR{\lif}}
\UnaryInf$ \Gamma \fCenter \Delta, A \lif B $
\DisplayProof\]
suggests we should take $\Gamma = \{C \lif (D \lif E)\}$, $\Delta =
\emptyset$, $A \ident C$ and $B \ident D \lif E$. So the last
inference of a proof could be
\[\Axiom$ C, (C \land D) \lif E \fCenter D \lif E$
\RightLabel{\RightR{\lif}}
\UnaryInf$(C \land D) \lif E \fCenter C \lif (D \lif E)$
\DisplayProof\]
If we repeat this idea, focusing on $D \lif E$ as active formula, we
generate
\[\Axiom$ D, C, (C \land D) \lif E \fCenter E$
\RightLabel{\RightR{\lif}}
\UnaryInf$C, (C \land D) \lif E \fCenter D \lif E$
\RightLabel{\RightR{\lif}}
\UnaryInf$(C \land D) \lif E \fCenter C \lif (D \lif E)$
\DisplayProof\]
Now we have to focus on $(C \land D) \lif E$ and use the $\LeftR{\lif}$ rule of \Log{G1c},
\[\Axiom$ \Gamma \fCenter \Delta, A$
\Axiom$ B, \Gamma \fCenter \Delta$
\RightLabel{\LeftR{\lif}}
\BinaryInf$ A \lif B, \Gamma \fCenter \Delta$
\DisplayProof\]
This leads us to
\[\Axiom$D, C \fCenter E, C \land D$
\Axiom$D, C, E \fCenter E$
\RightLabel{\LeftR{\lif}}
\BinaryInf$ D, C, (C \land D) \lif E \fCenter E$
\RightLabel{\RightR{\lif}}
\UnaryInf$C, (C \land D) \lif E \fCenter D \lif E$
\RightLabel{\RightR{\lif}}
\UnaryInf$(C \land D) \lif E \fCenter C \lif (D \lif E)$
\DisplayProof\]
Using the \olup{formula} $C \land D$ in the left topmost sequent as
principal and the $\RightR{\land}$ rule gives us
\[
\Axiom$D, C \fCenter E, C$
\Axiom$D, C \fCenter E, D$
\RightLabel{\RightR{\land}}
\BinaryInf$D, C \fCenter E, C \land D$
\Axiom$D, C, E \fCenter E$
\RightLabel{\RightR{\lif}}
\BinaryInf$ D, C, (C \land D) \lif E \fCenter E$
\RightLabel{\RightR{\lif}}
\UnaryInf$C, (C \land D) \lif E \fCenter D \lif E$
\RightLabel{\RightR{\lif}}
\UnaryInf$(C \land D) \lif E \fCenter C \lif (D \lif E)$
\DisplayProof
\]
So far so good. Note that everything we did also works in \Log{G3c},
since the rules we've made use of so far are the same in \Log{G1c}
and~\Log{G3c}. We've arrived at \olup{proof} fragment where all top
sequents contain the same \olup{formula} on the left as on the right. But
they are not quite axioms yet.

In \Log{G1c}, we have to give \olup{proof}s of the topmost sequents from
axioms of the form $A \Sequent A$. This is easily done using the
weakening rules, e.g., the leftmost topmost sequent $D, C \Sequent
E, C$ can be \olup{prove}d
as follows:
\[
\Axiom$C \fCenter C$
\RightLabel{\LeftR{\Weakening}}
\UnaryInf$D, C \fCenter C$
\RightLabel{\RightR{\Weakening}}
\UnaryInf$D, C \fCenter E, C$
\]
Overall, we have \olup{proof} in \Log{G1c} which looks like this:
\[
\Axiom$C \fCenter C$
\RightLabel{\LeftR{\Weakening}}
\UnaryInf$D, C \fCenter C$
\RightLabel{\RightR{\Weakening}}
\UnaryInf$D, C \fCenter E, C$
\Axiom$D \fCenter D$
\RightLabel{\LeftR{\Weakening}}
\UnaryInf$D, C \fCenter D$
\RightLabel{\RightR{\Weakening}}
\UnaryInf$D, C \fCenter E, D$
\RightLabel{\RightR{\land}}
\BinaryInf$D, C \fCenter E, C \land D$
\Axiom$E \fCenter E$
\RightLabel{\LeftR{\Weakening}}
\UnaryInf$C, E \fCenter C$
\RightLabel{\LeftR{\Weakening}}
\UnaryInf$D, C, E \fCenter E$
\RightLabel{\LeftR{\lif}}
\BinaryInf$ D, C, (C \land D) \lif E \fCenter E$
\RightLabel{\RightR{\lif}}
\UnaryInf$C, (C \land D) \lif E \fCenter D \lif E$
\RightLabel{\RightR{\lif}}
\UnaryInf$(C \land D) \lif E \fCenter C \lif (D \lif E)$
\DisplayProof
\]
The structure of this proof (in particular, its height) is independent
of what the \olup{formula}s $C$, $D$, and~$E$ are. We can replace
these symbols in the \olup{proof} above by any \olup{formula}s and the
result will still be a \olup{proof} in~\Log{G1c}. The proof in \Log{G1c}
is \emph{schematic}.

In \Log{G3c}, axioms are sequents of the form $A, \Gamma \Sequent
\Delta, A$ where $A$ must be atomic. So although $D, C \Sequent
E, C$ looks like an axiom of \Log{G3c} ($\Gamma = \{D\}$; $\Delta =
\{E\}$), it is only \emph{really} an axiom if $C$ in fact is atomic.
Our proof fragment is only a complete proof in \Log{G3c} if $C$,
$D$, and $E$ are atomic \olup{formula}s.

Although we haven't, strictly speaking, shown that
\[\Log{G1c} \Proves (C \land D) \lif E \fCenter C \lif (D \lif
E),\] in practice this is all we need to do (and can do). That's
because every sequent of the form $A, \Gamma \Sequent \Delta, A$ is
provable in \Log{G3c}, even if $A$ itself is not atomic. We can show
this by induction on the depth~$\depth{A}$ of~$A$.

\begin{defn}\ollabel{defn:depth}
The \emph{depth}~$\depth{A}$ of \olup{formula} is defined as follows:
\begin{enumerate}
  \item If $A$ is atomic, $\depth{A} = 0$.
  \item If $A \ident \lnot B$, $A \ident \lforall[x][A]$, or 
  $A \ident \lexists[x][A]$, then $\depth{A} = \depth{B} + 1$.
  \item If $A \ident (B \land C)$, $(B \lor C)$, or $(B \lif
  C)$, then $\depth{A} = \max(\depth{B},\depth{C}) + 1$.
\end{enumerate}
\end{defn}

It is not hard to prove that $\depth{A(x)} = \depth{A(t)}$ for any
term~$t$. The important fact for us is that in any rule, the depth of
the principal \olup{formula} is always greater than the depth of the
active \olup{formula}s. For instance, we have:
\begin{align*}
\depth{P(x)} = \depth{P(a)} & = 0,\\
\depth{\lforall[x][P(x)]} & = 1,\\
\depth{(\lforall[x][P(x)] \lif P(a))} & = \max(1,0) + 1 = 2.
\end{align*}
We can use this to prove some results by induction on~$\depth{A}$,
such as the following.

\begin{prop}\ollabel{prop:G1c-axioms}
For any~$A$, there is \olup{proof} in $\Log{G1c}$ of $A \Sequent A$
in which all axioms are atomic.
\end{prop}

\begin{proof}
  By induction on~$\depth{A}$. If $\depth{A} = 0$ then $A$ is
  atomic and $A \Sequent A$ already is a \olup{proof} in~\Log{G1c} of
  the required form. If $\depth{A} > 0$, $A$ is constructed from
  \olup{formula}s of lower depth, e.g, $A \ident \lnot B$, $A \ident
  (B \land C)$, or $A \ident \lforall[x][A(x)]$. The inductive
  hypothesis is that for all $D$ with $\depth{D} < \depth{A}$,
  $\Log{G1c} \Proves D \Sequent D$ from atomic axioms.

  If $A \ident \lnot B$, then $\depth{B} < \depth{A}$ and by
  inductive hypothesis, there is \olup{proof}~$\pi_1$ in \Log{G1c} from
  atomic axioms of
  $B \Sequent B$. We can extend this to \olup{proof}
  of $\lnot B \Sequent \lnot B$ as follows:
  \[
  \AxiomC{}
  \RightLabel{$\pi_1$}
  \Deduce$B \fCenter B$
  \RightLabel{\LeftR{\lnot}}
  \UnaryInf$\lnot B, B \fCenter $
  \RightLabel{\RightR{\lnot}}
  \UnaryInf$\lnot B \fCenter \lnot B$
  \DisplayProof
\]

  If $A \ident (B \land C)$, then $\depth{B} < \depth{A}$ and
  $\depth{C} < \depth{A}$. By inductive hypothesis, there are
  \olup{proof}s~$\pi_1$ and~$\pi_2$ in \Log{G1c} from atomic axioms of $B
  \Sequent B$ and $C \Sequent C$. We can extend these
  to \olup{proof} of $B \land C \Sequent B \land C$
  as follows:
  \[
  \AxiomC{}
  \RightLabel{$\pi_1$}
  \Deduce$B, \fCenter B$
  \RightLabel{\LeftR{\Weakening}}
  \UnaryInf$B, C, \fCenter B$
  \RightLabel{\LeftR{\land}}
  \UnaryInf$B \land C, C\fCenter B$
  \RightLabel{\LeftR{\land}}
  \UnaryInf$B \land C, B \land C \fCenter B$
  \AxiomC{}
  \RightLabel{$\pi_2$}
  \Deduce$C, \fCenter C$
  \RightLabel{\LeftR{\Weakening}}
  \UnaryInf$B, C, \fCenter C$
  \RightLabel{\LeftR{\land}}
  \UnaryInf$B \land C, C \fCenter C$
  \RightLabel{\LeftR{\land}}
  \UnaryInf$B \land C, B \land C \fCenter C$
  \RightLabel{\RightR{\land}}
  \BinaryInf$B \land C \fCenter B \land C$
  \RightLabel{\LeftR{\Contraction}}
  \UnaryInf$B \land C \fCenter B$
  \DisplayProof
  \]

  Now suppose $A \ident \lforall[x][B(x)]$. Take \olup{constant}~$c$
  that does not occur in $\lforall[x][B(x)]$. Then $\depth{B(c)} =
  \depth{B(x)} < \depth{A}$ and by inductive hypothesis, there is
  \olup{proof}~$\pi_1$ in \Log{G1c} from atomic axioms of $B(c)
  \Sequent B(c)$. We can extend this to \olup{proof} of
  $\lforall[x][B(x)] \Sequent \lforall[x][B(x)]$ as
  follows:
  \[
  \AxiomC{}
  \RightLabel{$\pi_1$}
  \Deduce$B(c) \fCenter B(c)$
  \RightLabel{\LeftR{\lforall}}
  \UnaryInf$\lforall[x][B(x)] \fCenter B(c)$
  \RightLabel{\RightR{\lforall}}
  \UnaryInf$\lforall[x][B(x)] \fCenter \lforall[x][B(x)]$
  \DisplayProof
\]

  The remaining cases are left as exercises.
\end{proof}

\begin{prob}
  Complete the proof of \olref{prop:G1c-axioms} by carrying out the
  cases where $A \ident (B \lor C)$, $A \ident (B \lif C)$ and
  $A \ident \lexists[x][B(x)]$.
\end{prob}

For our purposes, we could also have used the \olup{proof}
\[
\AxiomC{}
\RightLabel{$\pi_1$}
\Deduce$B \fCenter B$
\RightLabel{\RightR{\lnot}}
\UnaryInf$\fCenter B, \lnot B$
\RightLabel{\LeftR{\lnot}}
\UnaryInf$\lnot B \fCenter \lnot B$
\DisplayProof
\]
in the case where $A \ident \lnot B$. However, there are sequent
systems in which this wouldn't have worked. The intuitionistic
sequent calculus \Log{G1i} is like \Log{G1c}, except it restricts all
sequents to at most one \olup{formula} on the right side of a sequent. If
$\Delta = \emptyset$, then the first \olup{proof} would also be
\olup{proof} in~\Log{G1i}, but the second would not be, since
one sequent contains both $B$ and~$\lnot B$ on the right.

Note that even in \Log{G1c} we could not have used an alternate order
of rules in the case where $A \ident \lforall[x][B(x)]$. The
following is \emph{not} \olup{proof}
  \[
  \AxiomC{}
  \RightLabel{$\pi_1$}
  \Deduce$B(c) \fCenter B(c)$
  \RightLabel{\RightR{\lforall}}
  \UnaryInf$B(c) \fCenter \lforall[x][B(x)]$
  \RightLabel{\LeftR{\lforall}}
  \UnaryInf$\lforall[x][B(x)] \fCenter \lforall[x][B(x)]$
  \DisplayProof
\]
since the eigenvariable condition on~$\RightR{\lforall}$ is violated.

\begin{prop}\ollabel{prop:G3c-axioms}
Any sequent of the form $A, \Gamma
\Sequent \Delta, A$ (with $A$ \emph{not} necessarily atomic) is
provable in~\Log{G3c}.
\end{prop}

\begin{proof}
The proof is very similar to the one of \olref{prop:G1c-axioms}, but
we must be careful with the inductive hypothesis. The case where $n =
\depth{A} = 0$ is again easy. If $n=0$, $A$ is atomic, and $A,
\Gamma \Sequent \Delta, A$ is an axiom of~\Log{G3c}.

Now suppose $n>0$. Again we distinguish cases. The inductive
hypothesis is: for all $D!$ with $\depth{D} < n$, and all $\Pi$
and $\Lambda$, $\Log{G3c} \Proves D, \Pi \Sequent \Lambda, D$. Here
is the case for $A \ident (B \land C)$. We use the inductive
hypothesis twice: once for $D = B$, $\Pi = C, \Gamma$, and $\Lambda
= \Delta$ to get \olup{proof}~$\pi_1$ of $B, C, \Gamma \Sequent
\Delta, B$, and once for $D = C$, $\Pi = B, \Lambda$, and $\Lambda
= \Delta$ to get \olup{proof}~$\pi_2$ of $B, C, \Gamma \Sequent
\Delta, C$. Here is the \olup{proof} of~$B \land C, \Gamma \Sequent
\Delta, B \land C$:
\[
\AxiomC{}
\RightLabel{$\pi_1$}
\Deduce$B, C, \Gamma \fCenter \Delta, B$
\RightLabel{\LeftR{\land}}
\UnaryInf$B \land C, \Gamma \fCenter \Delta, B$
\AxiomC{}
\RightLabel{$\pi_2$}
\Deduce$B, C, \Gamma \fCenter \Delta, C$
\RightLabel{\LeftR{\land}}
\UnaryInf$B \land C, \Gamma \fCenter \Delta, C$
\RightLabel{\RightR{\land}}
\BinaryInf$B \land C, \Gamma \fCenter \Delta, B \land C$
\DisplayProof
\]

For $A \ident \lforall[x][B(x)]$, pick \olup{constant} not
in~$\lforall[x][B(x)]$, $\Gamma$, or~$\Delta$. Apply the inductive
hypothesis for $D = B(c)$, $\Pi = \lforall[x][B(x)], \Gamma$
and $\Lambda = \Delta$. This gives \olup{proof}~$\pi_1$ of $B(c),
\lforall[x][B(x)], \Gamma \Sequent \Delta, B(c)$, from which we
obtain:
\[
\AxiomC{}
\RightLabel{$\pi_1$}
\Deduce$B(c), \lforall[x][B(x)], \Gamma \fCenter \Delta, B(c)$
\RightLabel{\LeftR{\lforall}}
\UnaryInf$\lforall[x][B(x)], \Gamma \fCenter \Delta, B(c)$
\RightLabel{\RightR{\lforall}}
\UnaryInf$\lforall[x][B(x)], \Gamma \fCenter \Delta, \lforall[x][B(x)]$
\DisplayProof
\]
The eigenvariable condition on $\RightR{\lforall}$ is satisfied by the
way we chose~$c$.
\end{proof}

\begin{prob}
  Continue the proof of \olref{prop:G3c-axioms} by carrying out the
  cases where $A \ident (B \lif C)$ and $A \ident
  \lexists[x][B(x)]$.
\end{prob}

% ==== OLP file: content/proof-theory/sequent-calculus/quantifiers.tex ====
% Part: proof-theory
% Chapter: sequent-calculus
% Section: quantifiers



\olfileid{pt}{seq}{qua}

\olsection{Regular Proofs and Substitution}

The rules for the quantifiers introduce some complications owing to
the ``eigenvariable conditions'' that apply to the \LeftR{\lexists}
and~\RightR{\lforall} rules. Most of these complications can be dealt
with by ensuring that the \olup{constant}~$a$ in these inferences is
never reused. Let's introduce the following definition.

\begin{defn}
!!^a{proof}~$\pi$ in \Log{G1c} or \Log{G3c} is \emph{regular} if
\begin{enumerate}
  \item no eigenvariable~$a$ is used as the eigenvariable of more than
  one \LeftR{\lexists} or~\RightR{\lforall} inference, and
  \item if $a$ is the eigenvariable of a \LeftR{\lexists}
  or~\RightR{\lforall} inference, it only occurs above that inference
  in~$\pi$.
\end{enumerate}
\end{defn}

\begin{lem}\ollabel{lem:subst-regular} If $\pi$ is regular, ends in
$\Gamma \Sequent \Delta$, $c$ is \olup{constant} not used as an
eigenvariable in~$\pi$, and $t$~is a term not containing any
eigenvariables of~$\pi$, then $\Subst{\pi}{t}{c}$ resulting from~$\pi$
by replacing every occurrence of~$c$ by~$t$ is a regular \olup{proof}.
The end-sequent of $\Subst{\pi}{t}{c}$ is $\Subst{\Gamma}{t}{c}
\Sequent \Subst{\Delta}{t}{c}$, where $\Subst{\Gamma}{t}{c} =
\Setabs{A(t)}{A(c) \in \Gamma}$.
\end{lem}

\begin{proof}
  By induction on $\pheight{\pi}$. If $\pheight{\pi} = 0$, it consists
  simply of an axiom. Then $\Subst{\pi}{t}{c}$ is also an axiom, since
  any \olup{formula}~$A(c)$ in it turns into $A(t)$.

  If $\pheight{\pi} > 0$ we distinguish cases according to the last
  inference in~$\pi$. The cases for structural rules (in \Log{G1c})
  and logical rules for propositional connectives are easy and left as
  exercises. We consider the cases where $\pi$ ends in a quantifier
  inference.
  \begin{enumerate}
    \item If the last inference is $\RightR{\lforall}$, $\pi$ is of
    the form
    \[
    \AxiomC{}
    \RightLabel{$\pi'$}
    \Deduce$\Gamma(c) \fCenter \Delta'(c), A(a,c)$
    \RightLabel{\RightR{\lforall}}
    \UnaryInf$\Gamma(c) \fCenter \Delta'(c), \lforall[x][A(x,c)]$
    \DisplayProof
    \]
    By induction hypothesis, $\Subst{\pi'}{t}{c}$ is a regular
    \olup{proof} of $\Gamma(t) \fCenter \Delta'(t), A(a,t)$. Since $a$
    is an eigenvariable of~$\pi$, $a$~does not occur in~$t$.
    Therefore, the last inference in $\Subst{\pi}{t}{c}$,
    \[
    \AxiomC{}
    \RightLabel{$\Subst{\pi'}{t}{c}$}
    \Deduce$\Gamma(t) \fCenter \Delta'(t), A(a,t)$
    \RightLabel{\RightR{\lforall}}
    \UnaryInf$\Gamma(t) \fCenter \Delta'(t), \lforall[x][A(x,t)]$
    \DisplayProof
    \]
    is correct: the conclusion does not contain~$a$. 
    \item If the last inference is $\LeftR{\lforall}$ in \Log{G3c},
    $\pi$ is of the form
    \[
    \AxiomC{}
    \RightLabel{$\pi'$}
    \Deduce$A(s(c),c), \lforall[x][A(x,c)], \Gamma(c) \fCenter \Delta'(c)$
    \RightLabel{\RightR{\lforall}}
    \UnaryInf$\lforall[x][A(x,c)], \Gamma(c) \fCenter \Delta'(c)$
    \DisplayProof
    \]
    By induction hypothesis, $\Subst{\pi'}{t}{c}$ is a regular
    \olup{proof} of $A(s(t),t), \lforall[x][A(x,t)], \Gamma(t) \fCenter
    \Delta'(t)$. The last inference in
    $\Subst{\pi}{t}{c}$,
    \[
    \AxiomC{}
    \RightLabel{$\Subst{\pi'}{t}{c}$}
    \Deduce$A(s(t),t), \lforall[x][A(x,t)], \Gamma(t) \fCenter \Delta'(t)$
    \RightLabel{\RightR{\lforall}}
    \UnaryInf$\lforall[x][A(x,t)], \Gamma(t) \fCenter \Delta'(t)$
    \DisplayProof
    \]
    is correct. The case for \LeftR{\lforall} in~\Log{G1c} is similar,
    but slightly less complicated.
    \item The last inference is \RightR{\lexists} or \LeftR{\lexists}: exercise.
  \end{enumerate}
  In each case, we added a sequent to the end of a regular \olup{proof}
  (or two regular \olup{proof}s in the case of two-premise rules). The
  new sequent differs from the end-sequent of $\pi$ only in that $c$
  was replaced by the term~$t$. Moreover, the eigenvariables of $\pi$
  and Since $t$ was assumed not to contain an eigenvariable of~$\pi$
  and $\Subst{\pi}{t}{c}$ are the same, the new end-sequent does not
  contain an eigenvariable of $\Subst{\pi}{t}{c}$. Consequently,
  $\Subst{\pi}{t}{c}$ is regular.
\end{proof}

\begin{prop}
If $\pi$ is \olup{proof} in \Log{G1c} or \Log{G3c}, then there is a
regular \olup{proof}~$\pi'$ with the same structure (in particular, the
same \olup{height}) which differs from $\pi$ only by replacing
eigenvariables.
\end{prop}

\begin{proof}
For the purpose of the proof only, call an eigenvariable inference
of~$\pi$ \emph{clean} if its eigenvariable is not the eigenvariable of
another inference nor does it occur anywhere in~$\pi$ except the
immediate sub-\olup{proof} of the inference, and \emph{dirty} otherwise.
Let $n$ be the number of dirty eigenvariable inferences in~$\pi$. We
show, by induction on~$n$, that $\pi$~can be transformed into the
required regular \olup{proof}~$\pi'$. If $n=0$, all eigenvariable
inferences in~$\pi$ are clean, so and $\pi$ is regular, and there is
nothing to prove. Otherwise, pick a highest eigenvariable inference,
say \RightR{\lforall}. The sub-\olup{proof}~$\pi_1$ of~$\pi$ ending in it
is of the form
    \[
    \AxiomC{}
    \RightLabel{$\pi_2$}
    \Deduce$\Gamma \fCenter \Delta, A(c)$
    \RightLabel{\RightR{\lforall}}
    \UnaryInf$\Gamma \fCenter \Delta, \lforall[x][A(x)]$
    \DisplayProof
    \]
Note that as $c$ is an eigenvariable, it does not occur in~$\Gamma$,
$\Delta$, or $\lforall[x][A(x)]$. The proof $\pi_2$ is regular, since
it contains no eigenvariable inferences at all. Let $a$ be
\olup{constant} not occuring in~$\pi$. By \olref{lem:subst-regular},
$\Subst{\pi_2}{a}{c}$ is a regular proof of $\Gamma \fCenter \Delta,
A(a)$, to which we can apply~$\RightR{\lforall}$. If we replace $\pi_1$ in~$\pi$ by the proof $\pi_1'$,
    \[
    \AxiomC{}
    \RightLabel{$\Subst{\pi_2}{a}{c}$}
    \Deduce$\Gamma \fCenter \Delta, A(a)$
    \RightLabel{\RightR{\lforall}}
    \UnaryInf$\Gamma \fCenter \Delta, \lforall[x][A(x)]$
    \DisplayProof
    \]
we have replaced a dirty eigenvariable inference by a clean one, and
the only difference is the replacement of the eigenvariable $c$
by~$a$. The resulting proof has $n-1$ dirty eigenvariable inferences,
and the result follows from the inductive hypothesis.
\end{proof}

We will not appeal to the proposition just proved explicitly, but we
will assume that all \olup{proof}s considered are regular---if they are
not, we can always replace eigenvariables to make them regular.

% ==== OLP file: content/proof-theory/sequent-calculus/rules-G1c.tex ====
% Part: proof-theory
% Chapter: sequent-calculus
% Section: rules-G1c



\begin{table}
\begin{tabular}{@{}c@{}c@{}}
\multicolumn{2}{@{}c@{}}{Axioms}\\[2ex]
$A \Sequent A$ & $\lfalse \Sequent \quad$
\\[2ex]
\multicolumn{2}{@{}c@{}}{Structural Rules}\\[2ex]
\Axiom$ \Gamma \fCenter \Delta $
\RightLabel{\LeftR{\Weakening}}
\UnaryInf$ A, \Gamma \fCenter \Delta$
\DisplayProof
&
\Axiom$ \Gamma \fCenter \Delta$
\RightLabel{\RightR{\Weakening}}
\UnaryInf$ \Gamma \fCenter \Delta, A$
\DisplayProof
\\[4ex]
\Axiom$ A, A, \Gamma \fCenter \Delta $
\RightLabel{\LeftR{\Contraction}}
\UnaryInf$ A, \Gamma \fCenter \Delta$
\DisplayProof
&
\Axiom$ \Gamma \fCenter \Delta, A, A$
\RightLabel{\RightR{\Contraction}}
\UnaryInf$ \Gamma \fCenter \Delta, A$
\DisplayProof
\\[4ex]
\multicolumn{2}{@{}c@{}}{Logical Rules}\\[2ex]
\Axiom$ \Gamma \fCenter \Delta, A $
\RightLabel{\LeftR{\lnot}}
\UnaryInf$ \lnot A, \Gamma \fCenter \Delta$
\DisplayProof
&
\Axiom$A, \Gamma \fCenter \Delta$
\RightLabel{\RightR{\lnot}}
\UnaryInf$ \Gamma \fCenter \Delta, \lnot A $
\DisplayProof
\\[4ex]\begin{tabular}[t]{@{}l@{}}
\Axiom$ A, \Gamma \fCenter \Delta$
\RightLabel{\LeftR{\land}}
\UnaryInf$ A \land B, \Gamma \fCenter \Delta$
\DisplayProof
\\[3ex]
\Axiom$B, \Gamma \fCenter \Delta$
\RightLabel{\LeftR{\land}}
\UnaryInf$A \land B, \Gamma \fCenter \Delta$
\DisplayProof\\[3ex]
\end{tabular}
&
\Axiom$\Gamma \fCenter \Delta, A$
\Axiom$ \Gamma \fCenter \Delta, B$
\RightLabel{\RightR{\land}}
\BinaryInf$ \Gamma \fCenter \Delta, A \land B $
\DisplayProof
\\[4ex]\Axiom$A, \Gamma \fCenter \Delta$
\Axiom$B, \Gamma \fCenter \Delta$
\RightLabel{\LeftR{\lor}}
\BinaryInf$A \lor B, \Gamma \fCenter \Delta$
\DisplayProof
&
\begin{tabular}[t]{@{}r@{}}
\Axiom$\Gamma \fCenter \Delta, A$
\RightLabel{\RightR{\lor}}
\UnaryInf$ \Gamma \fCenter \Delta, A \lor B$
\DisplayProof
\\[3ex]
\Axiom$ \Gamma \fCenter \Delta, B$
\RightLabel{\RightR{\lor}}
\UnaryInf$ \Gamma \fCenter \Delta, A \lor B$
\DisplayProof\\[3ex]
\end{tabular}
\\[4ex]
\Axiom$ \Gamma \fCenter \Delta, A$
\Axiom$ B, \Gamma \fCenter \Delta$
\RightLabel{\LeftR{\lif}}
\BinaryInf$ A \lif B, \Gamma \fCenter \Delta$
\DisplayProof
&
\Axiom$ A, \Gamma \fCenter \Delta, B$
\RightLabel{\RightR{\lif}}
\UnaryInf$ \Gamma \fCenter \Delta, A \lif B $
\DisplayProof
\\[4ex]
\Axiom$ A(t), \Gamma \fCenter \Delta$
\RightLabel{\LeftR{\lforall}}
\UnaryInf$ \lforall[x][A(x)],\Gamma \fCenter \Delta$
\DisplayProof
&
\Axiom$ \Gamma \fCenter \Delta, A(a) $
\RightLabel{\RightR{\lforall}}
\UnaryInf$ \Gamma \fCenter \Delta, \lforall[x][A(x)]$
\DisplayProof
\\[4ex]
\Axiom$ A(a), \Gamma \fCenter \Delta $
\RightLabel{\LeftR{\lexists}}
\UnaryInf$ \lexists[x][A(x)], \Gamma \fCenter \Delta$
\DisplayProof
&
\Axiom$ \Gamma \fCenter \Delta, A(t) $
\RightLabel{\RightR{\lexists}}
\UnaryInf$ \Gamma \fCenter \Delta, \lexists[x][A(x)]$
\DisplayProof
\end{tabular}
\caption{Rules of \Log{G1c}. In $\RightR{\forall}$ and
$\LeftR{\exists}$, the \olup{constant}~$a$ must not occur in the
conclusion sequent. $\Gamma$~and $\Delta$ are
multisets of \olup{formula}s.}
\ollabel{tab:G1c}
\end{table}

% ==== OLP file: content/proof-theory/sequent-calculus/rules-G1i.tex ====
% Part: proof-theory
% Chapter: sequent-calculus
% Section: rules-G1c



\begin{table}
\begin{tabular}{@{}c@{}c@{}}
\multicolumn{2}{@{}c@{}}{Axioms}\\[2ex]
$A \Sequent A$ & $\lfalse \Sequent \quad$
\\[2ex]
\multicolumn{2}{@{}c@{}}{Structural Rules}\\[2ex]
\Axiom$ \Gamma \fCenter \Delta $
\RightLabel{\LeftR{\Weakening}}
\UnaryInf$ A, \Gamma \fCenter \Delta$
\DisplayProof
&
\Axiom$ \Gamma \fCenter $
\RightLabel{\RightR{\Weakening}}
\UnaryInf$ \Gamma \fCenter A$
\DisplayProof
\\[4ex]
\Axiom$ A, A, \Gamma \fCenter \Delta $
\RightLabel{\LeftR{\Contraction}}
\UnaryInf$ A, \Gamma \fCenter \Delta$
\DisplayProof
&
\\[4ex]
\multicolumn{2}{@{}c@{}}{Logical Rules}\\[2ex]
\Axiom$ \Gamma \fCenter A $
\RightLabel{\LeftR{\lnot}}
\UnaryInf$ \lnot A, \Gamma \fCenter $
\DisplayProof
&
\Axiom$A, \Gamma \fCenter $
\RightLabel{\RightR{\lnot}}
\UnaryInf$ \Gamma \fCenter \lnot A $
\DisplayProof
\\[4ex]\begin{tabular}[t]{@{}l@{}}
\Axiom$ A, \Gamma \fCenter \Delta$
\RightLabel{\LeftR{\land}}
\UnaryInf$ A \land B, \Gamma \fCenter \Delta$
\DisplayProof
\\[3ex]
\Axiom$B, \Gamma \fCenter \Delta$
\RightLabel{\LeftR{\land}}
\UnaryInf$A \land B, \Gamma \fCenter \Delta$
\DisplayProof\\[3ex]
\end{tabular}
&
\Axiom$\Gamma \fCenter A$
\Axiom$ \Gamma \fCenter B$
\RightLabel{\RightR{\land}}
\BinaryInf$ \Gamma \fCenter A \land B $
\DisplayProof
\\[4ex]\Axiom$A, \Gamma \fCenter \Delta$
\Axiom$B, \Gamma \fCenter \Delta$
\RightLabel{\LeftR{\lor}}
\BinaryInf$A \lor B, \Gamma \fCenter \Delta$
\DisplayProof
&
\begin{tabular}[t]{@{}r@{}}
\Axiom$\Gamma \fCenter A$
\RightLabel{\RightR{\lor}}
\UnaryInf$ \Gamma \fCenter A \lor B$
\DisplayProof
\\[3ex]
\Axiom$ \Gamma \fCenter B$
\RightLabel{\RightR{\lor}}
\UnaryInf$ \Gamma \fCenter A \lor B$
\DisplayProof\\[3ex]
\end{tabular}
\\[4ex]
\Axiom$ \Gamma \fCenter A$
\Axiom$ B, \Gamma \fCenter \Delta$
\RightLabel{\LeftR{\lif}}
\BinaryInf$ A \lif B, \Gamma \fCenter \Delta$
\DisplayProof
&
\Axiom$ A, \Gamma \fCenter B$
\RightLabel{\RightR{\lif}}
\UnaryInf$ \Gamma \fCenter A \lif B $
\DisplayProof
\\[4ex]
\Axiom$ A(t), \Gamma \fCenter \Delta$
\RightLabel{\LeftR{\lforall}}
\UnaryInf$ \lforall[x][A(x)],\Gamma \fCenter \Delta$
\DisplayProof
&
\Axiom$ \Gamma \fCenter A(a) $
\RightLabel{\RightR{\lforall}}
\UnaryInf$ \Gamma \fCenter \lforall[x][A(x)]$
\DisplayProof
\\[4ex]
\Axiom$ A(a), \Gamma \fCenter \Delta $
\RightLabel{\LeftR{\lexists}}
\UnaryInf$ \lexists[x][A(x)], \Gamma \fCenter \Delta$
\DisplayProof
&
\Axiom$ \Gamma \fCenter A(t) $
\RightLabel{\RightR{\lexists}}
\UnaryInf$ \Gamma \fCenter \lexists[x][A(x)]$
\DisplayProof
\end{tabular}
\caption{Rules of \Log{G1i} for intuitionistic logic. In
$\RightR{\forall}$ and $\LeftR{\exists}$, the \olup{constant}~$a$ must
not occur in the conclusion sequent. $\Gamma$ is a multisets of
\olup{formula}s, and $\Delta$ may be empty or contain a single
\olup{formula}. There is no \RightR{\Contraction} rule. \Log{G1m} is
\Log{G1i} without~\RightR{\Weakening} and without axioms $\lfalse,
\Gamma \Sequent \Delta$.}
\ollabel{tab:G1i}
\end{table}

% ==== OLP file: content/proof-theory/sequent-calculus/rules-G2c.tex ====
% Part: proof-theory
% Chapter: sequent-calculus
% Section: rules-G1c



\begin{table}
\begin{tabular}{@{}c@{\hskip-5ex}c@{}}
\multicolumn{2}{@{}c@{}}{Axioms}\\[2ex]
$A \Sequent A$ & $\lfalse \Sequent \quad$
\\[2ex]
\multicolumn{2}{@{}c@{}}{Structural Rules}\\[2ex]
\Axiom$ \Gamma \fCenter \Delta $
\RightLabel{\LeftR{\Weakening}}
\UnaryInf$ A, \Gamma \fCenter \Delta$
\DisplayProof
&
\Axiom$ \Gamma \fCenter \Delta$
\RightLabel{\RightR{\Weakening}}
\UnaryInf$ \Gamma \fCenter \Delta, A$
\DisplayProof
\\[4ex]
\Axiom$ A, A, \Gamma \fCenter \Delta $
\RightLabel{\LeftR{\Contraction}}
\UnaryInf$ A, \Gamma \fCenter \Delta$
\DisplayProof
&
\Axiom$ \Gamma \fCenter \Delta, A, A$
\RightLabel{\RightR{\Contraction}}
\UnaryInf$ \Gamma \fCenter \Delta, A$
\DisplayProof
\\[4ex]
\multicolumn{2}{@{}c@{}}{Logical Rules}\\[2ex]
\Axiom$ \Gamma \fCenter \Delta, A $
\RightLabel{\LeftR{\lnot}}
\UnaryInf$ \lnot A, \Gamma \fCenter \Delta$
\DisplayProof
&
\Axiom$A, \Gamma \fCenter \Delta$
\RightLabel{\RightR{\lnot}}
\UnaryInf$ \Gamma \fCenter \Delta, \lnot A $
\DisplayProof
\\[4ex]\begin{tabular}[t]{@{}l@{}}
\Axiom$ A, \Gamma \fCenter \Delta$
\RightLabel{\LeftR{\land}}
\UnaryInf$ A \land B, \Gamma \fCenter \Delta$
\DisplayProof
\\[3ex]
\Axiom$B, \Gamma \fCenter \Delta$
\RightLabel{\LeftR{\land}}
\UnaryInf$A \land B, \Gamma \fCenter \Delta$
\DisplayProof\\[3ex]
\end{tabular}
&
\Axiom$\Gamma_1 \fCenter \Delta_1, A$
\Axiom$ \Gamma_2 \fCenter \Delta_2, B$
\RightLabel{\RightR{\land}}
\BinaryInf$ \Gamma_1, \Gamma_2 \fCenter \Delta_1, \Delta_2, A \land B $
\DisplayProof
\\[4ex]\Axiom$A, \Gamma_1 \fCenter \Delta_1$
\Axiom$B, \Gamma_2 \fCenter \Delta_2$
\RightLabel{\LeftR{\lor}}
\BinaryInf$A \lor B, \Gamma_1, \Gamma_2 \fCenter \Delta_1, \Delta_2$
\DisplayProof
&
\begin{tabular}[t]{@{}r@{}}
\Axiom$\Gamma \fCenter \Delta, A$
\RightLabel{\RightR{\lor}}
\UnaryInf$ \Gamma \fCenter \Delta, A \lor B$
\DisplayProof
\\[3ex]
\Axiom$ \Gamma \fCenter \Delta, B$
\RightLabel{\RightR{\lor}}
\UnaryInf$ \Gamma \fCenter \Delta, A \lor B$
\DisplayProof\\[3ex]
\end{tabular}
\\[4ex]
\Axiom$ \Gamma_1 \fCenter \Delta_1, A$
\Axiom$ B, \Gamma_2 \fCenter \Delta_2$
\RightLabel{\LeftR{\lif}}
\BinaryInf$ A \lif B, \Gamma_1, \Gamma_2 \fCenter \Delta_1, \Delta_2$
\DisplayProof
&
\Axiom$ A, \Gamma \fCenter \Delta, B$
\RightLabel{\RightR{\lif}}
\UnaryInf$ \Gamma \fCenter \Delta, A \lif B $
\DisplayProof
\\[4ex]
\Axiom$ A(t), \Gamma \fCenter \Delta$
\RightLabel{\LeftR{\lforall}}
\UnaryInf$ \lforall[x][A(x)],\Gamma \fCenter \Delta$
\DisplayProof
&
\Axiom$ \Gamma \fCenter \Delta, A(a) $
\RightLabel{\RightR{\lforall}}
\UnaryInf$ \Gamma \fCenter \Delta, \lforall[x][A(x)]$
\DisplayProof
\\[4ex]
\Axiom$ A(a), \Gamma \fCenter \Delta $
\RightLabel{\LeftR{\lexists}}
\UnaryInf$ \lexists[x][A(x)], \Gamma \fCenter \Delta$
\DisplayProof
&
\Axiom$ \Gamma \fCenter \Delta, A(t) $
\RightLabel{\RightR{\lexists}}
\UnaryInf$ \Gamma \fCenter \Delta, \lexists[x][A(x)]$
\DisplayProof
\end{tabular}
\caption{Rules of \Log{G2c}. In $\RightR{\forall}$ and
$\LeftR{\exists}$, the \olup{constant}~$a$ must not occur in the
conclusion sequent. $\Gamma$~and $\Delta$ are
multisets of \olup{formula}s.}
\ollabel{tab:G2c}
\end{table}

% ==== OLP file: content/proof-theory/sequent-calculus/rules-G3c.tex ====
% Part: proof-theory
% Chapter: sequent-calculus
% Section: rules-G3c



\begin{table}
\begin{tabular}{@{}c@{}c@{}}
\multicolumn{2}{@{}c@{}}{Axioms}\\[2ex]
$C, \Gamma \Sequent \Delta, C$ & $\lfalse, \Gamma \Sequent \Delta$
\\[2ex]
\multicolumn{2}{@{}c@{}}{Logical Rules}\\[2ex]
\Axiom$ \Gamma \fCenter \Delta, A $
\RightLabel{\LeftR{\lnot}}
\UnaryInf$ \lnot A, \Gamma \fCenter \Delta$
\DisplayProof
&
\Axiom$A, \Gamma \fCenter \Delta$
\RightLabel{\RightR{\lnot}}
\UnaryInf$ \Gamma \fCenter \Delta, \lnot A $
\DisplayProof
\\[4ex]
\Axiom$ A, B, \Gamma \fCenter \Delta$
\RightLabel{\LeftR{\land}}
\UnaryInf$ A \land B, \Gamma \fCenter \Delta$
\DisplayProof
&
\Axiom$\Gamma \fCenter \Delta, A$
\Axiom$ \Gamma \fCenter \Delta, B$
\RightLabel{\RightR{\land}}
\BinaryInf$ \Gamma \fCenter \Delta, A \land B $
\DisplayProof
\\[4ex]\Axiom$A, \Gamma \fCenter \Delta$
\Axiom$B, \Gamma \fCenter \Delta$
\RightLabel{\LeftR{\lor}}
\BinaryInf$A \lor B, \Gamma \fCenter \Delta$
\DisplayProof
&
\Axiom$\Gamma \fCenter \Delta, A, B$
\RightLabel{\RightR{\lor}}
\UnaryInf$ \Gamma \fCenter \Delta, A \lor B$
\DisplayProof
\\[4ex]
\Axiom$ \Gamma \fCenter \Delta, A$
\Axiom$ B, \Gamma \fCenter \Delta$
\RightLabel{\LeftR{\lif}}
\BinaryInf$ A \lif B, \Gamma \fCenter \Delta$
\DisplayProof
&
\Axiom$ A, \Gamma \fCenter \Delta, B$
\RightLabel{\RightR{\lif}}
\UnaryInf$ \Gamma \fCenter \Delta, A \lif B $
\DisplayProof
\\[4ex]
\Axiom$ A(t), \lforall[x][A(x)]\Gamma \fCenter \Delta$
\RightLabel{\LeftR{\lforall}}
\UnaryInf$ \lforall[x][A(x)],\Gamma \fCenter \Delta$
\DisplayProof
&
\Axiom$ \Gamma \fCenter \Delta, A(a) $
\RightLabel{\RightR{\lforall}}
\UnaryInf$ \Gamma \fCenter \Delta, \lforall[x][A(x)]$
\DisplayProof
\\[4ex]
\Axiom$ A(a), \Gamma \fCenter \Delta $
\RightLabel{\LeftR{\lexists}}
\UnaryInf$ \lexists[x][A(x)], \Gamma \fCenter \Delta$
\DisplayProof
&
\Axiom$ \Gamma \fCenter \Delta, \lexists[x][A(x)], A(t) $
\RightLabel{\RightR{\lexists}}
\UnaryInf$ \Gamma \fCenter \Delta, \lexists[x][A(x)]$
\DisplayProof
\end{tabular}
\caption{Rules of \Log{G3c}. In $\RightR{\forall}$ and
$\LeftR{\exists}$, the \olup{constant}~$a$ must not occur in the
conclusion sequent. $\Gamma$~and $\Delta$ are multisets of
\olup{formula}s. In axioms, the \olup{formula}~$C$ must be atomic.}
\ollabel{tab:G3c}
\end{table}

% ==== OLP file: content/proof-theory/sequent-calculus/rules-G3i.tex ====
% Part: proof-theory
% Chapter: sequent-calculus
% Section: rules-G3c



\begin{table}
\begin{tabular}{@{}c@{}c@{}}
\multicolumn{2}{@{}c@{}}{Axioms}\\[2ex]
$C, \Gamma \Sequent C$ & $\lfalse, \Gamma \Sequent \Delta$
\\[2ex]
\multicolumn{2}{@{}c@{}}{Logical Rules}\\[2ex]
\Axiom$\lnot A, \Gamma \fCenter A $
\RightLabel{\LeftR{\lnot}}
\UnaryInf$\lnot A, \Gamma \fCenter \Delta$
\DisplayProof
&
\Axiom$A, \Gamma \fCenter$
\RightLabel{\RightR{\lnot}}
\UnaryInf$ \Gamma \fCenter \lnot A $
\DisplayProof
\\[4ex]
\Axiom$ A, B, \Gamma \fCenter \Delta$
\RightLabel{\LeftR{\land}}
\UnaryInf$ A \land B, \Gamma \fCenter \Delta$
\DisplayProof
&
\Axiom$\Gamma \fCenter A$
\Axiom$ \Gamma \fCenter B$
\RightLabel{\RightR{\land}}
\BinaryInf$ \Gamma \fCenter A \land B $
\DisplayProof
\\[4ex]\Axiom$A, \Gamma \fCenter \Delta$
\Axiom$B, \Gamma \fCenter \Delta$
\RightLabel{\LeftR{\lor}}
\BinaryInf$A \lor B, \Gamma \fCenter \Delta$
\DisplayProof
&
\begin{tabular}[t]{@{}r@{}}
\Axiom$\Gamma \fCenter \Delta, A$
\RightLabel{\RightR{\lor}}
\UnaryInf$ \Gamma \fCenter \Delta, A \lor B$
\DisplayProof
\\[3ex]
\Axiom$ \Gamma \fCenter \Delta, B$
\RightLabel{\RightR{\lor}}
\UnaryInf$ \Gamma \fCenter \Delta, A \lor B$
\DisplayProof\\[3ex]
\end{tabular}
\\[4ex]
\Axiom$A \lif B, \Gamma \fCenter A$
\Axiom$B, \Gamma \fCenter \Delta$
\RightLabel{\LeftR{\lif}}
\BinaryInf$A \lif B, \Gamma \fCenter \Delta$
\DisplayProof
&
\Axiom$A, \Gamma \fCenter B$
\RightLabel{\RightR{\lif}}
\UnaryInf$ \Gamma \fCenter A \lif B $
\DisplayProof
\\[4ex]
\Axiom$ A(t), \lforall[x][A(x)]\Gamma \fCenter \Delta$
\RightLabel{\LeftR{\lforall}}
\UnaryInf$ \lforall[x][A(x)],\Gamma \fCenter \Delta$
\DisplayProof
&
\Axiom$ \Gamma \fCenter A(a) $
\RightLabel{\RightR{\lforall}}
\UnaryInf$ \Gamma \fCenter \lforall[x][A(x)]$
\DisplayProof
\\[4ex]
\Axiom$ A(a), \Gamma \fCenter \Delta $
\RightLabel{\LeftR{\lexists}}
\UnaryInf$ \lexists[x][A(x)], \Gamma \fCenter \Delta$
\DisplayProof
&
\Axiom$ \Gamma \fCenter A(t) $
\RightLabel{\RightR{\lexists}}
\UnaryInf$ \Gamma \fCenter \lexists[x][A(x)]$
\DisplayProof
\end{tabular}
\caption{Rules of \Log{G3i} for intuitionistic logic. In
$\RightR{\forall}$ and $\LeftR{\exists}$, the \olup{constant}~$a$ must
not occur in the conclusion sequent. $\Gamma$ is a multisets of
\olup{formula}s, and $\Delta$ may be empty or contain a single
\olup{formula}. In axioms, the \olup{formula}~$C$ must be atomic.
\Log{G1m} is \Log{G1i} without axioms $\lfalse, \Gamma \Sequent
\Delta$.}
\ollabel{tab:G3i}
\end{table}

% ==== OLP file: content/proof-theory/sequent-calculus/rules-LK.tex ====
% Part: proof-theory
% Chapter: sequent-calculus
% Section: rules-LK



\begin{table}
\begin{tabular}{@{}c@{}c@{}}
\multicolumn{2}{@{}c@{}}{Axioms}\\[2ex]
$A \Sequent A$ & $\lfalse \Sequent \quad$
\\[2ex]
\multicolumn{2}{@{}c@{}}{Structural Rules}\\[2ex]
\Axiom$ \Gamma \fCenter \Delta $
\RightLabel{\LeftR{\Weakening}}
\UnaryInf$ A, \Gamma \fCenter \Delta$
\DisplayProof
&
\Axiom$ \Gamma \fCenter \Delta$
\RightLabel{\RightR{\Weakening}}
\UnaryInf$ \Gamma \fCenter \Delta, A$
\DisplayProof
\\[4ex]
\Axiom$ A, A, \Gamma \fCenter \Delta $
\RightLabel{\LeftR{\Contraction}}
\UnaryInf$ A, \Gamma \fCenter \Delta$
\DisplayProof
&
\Axiom$ \Gamma \fCenter \Delta, A, A$
\RightLabel{\RightR{\Contraction}}
\UnaryInf$ \Gamma \fCenter \Delta, A$
\DisplayProof
\\[4ex]
\Axiom$ \Gamma, A, B, \Pi \fCenter \Delta $
\RightLabel{\LeftR{\Exchange}}
\UnaryInf$\Gamma, B, A, \Pi \fCenter \Delta$
\DisplayProof
&
\Axiom$ \Gamma \fCenter \Delta, A, B, \Lambda$
\RightLabel{\RightR{\Exchange}}
\UnaryInf$ \Gamma \fCenter \Delta, B, A, \Lambda$
\DisplayProof
\\[4ex]
\multicolumn{2}{@{}c@{}}{Logical Rules}\\[2ex]
\Axiom$ \Gamma \fCenter \Delta, A $
\RightLabel{\LeftR{\lnot}}
\UnaryInf$ \lnot A, \Gamma \fCenter \Delta$
\DisplayProof
&
\Axiom$A, \Gamma \fCenter \Delta$
\RightLabel{\RightR{\lnot}}
\UnaryInf$ \Gamma \fCenter \Delta, \lnot A $
\DisplayProof
\\[4ex]\begin{tabular}[t]{@{}l@{}}
\Axiom$ A, \Gamma \fCenter \Delta$
\RightLabel{\LeftR{\land}}
\UnaryInf$ A \land B, \Gamma \fCenter \Delta$
\DisplayProof
\\[3ex]
\Axiom$B, \Gamma \fCenter \Delta$
\RightLabel{\LeftR{\land}}
\UnaryInf$A \land B, \Gamma \fCenter \Delta$
\DisplayProof\\[3ex]
\end{tabular}
&
\Axiom$\Gamma \fCenter \Delta, A$
\Axiom$ \Gamma \fCenter \Delta, B$
\RightLabel{\RightR{\land}}
\BinaryInf$ \Gamma \fCenter \Delta, A \land B $
\DisplayProof
\\[4ex]\Axiom$A, \Gamma \fCenter \Delta$
\Axiom$B, \Gamma \fCenter \Delta$
\RightLabel{\LeftR{\lor}}
\BinaryInf$A \lor B, \Gamma \fCenter \Delta$
\DisplayProof
&
\begin{tabular}[t]{@{}r@{}}
\Axiom$\Gamma \fCenter \Delta, A$
\RightLabel{\RightR{\lor}}
\UnaryInf$ \Gamma \fCenter \Delta, A \lor B$
\DisplayProof
\\[3ex]
\Axiom$ \Gamma \fCenter \Delta, B$
\RightLabel{\RightR{\lor}}
\UnaryInf$ \Gamma \fCenter \Delta, A \lor B$
\DisplayProof\\[3ex]
\end{tabular}
\\[4ex]
\Axiom$ \Gamma \fCenter \Delta, A$
\Axiom$ B, \Gamma \fCenter \Delta$
\RightLabel{\LeftR{\lif}}
\BinaryInf$ A \lif B, \Gamma \fCenter \Delta$
\DisplayProof
&
\Axiom$ A, \Gamma \fCenter \Delta, B$
\RightLabel{\RightR{\lif}}
\UnaryInf$ \Gamma \fCenter \Delta, A \lif B $
\DisplayProof
\\[4ex]
\Axiom$ A(t), \Gamma \fCenter \Delta$
\RightLabel{\LeftR{\lforall}}
\UnaryInf$ \lforall[x][A(x)],\Gamma \fCenter \Delta$
\DisplayProof
&
\Axiom$ \Gamma \fCenter \Delta, A(a) $
\RightLabel{\RightR{\lforall}}
\UnaryInf$ \Gamma \fCenter \Delta, \lforall[x][A(x)]$
\DisplayProof
\\[4ex]
\Axiom$ A(a), \Gamma \fCenter \Delta $
\RightLabel{\LeftR{\lexists}}
\UnaryInf$ \lexists[x][A(x)], \Gamma \fCenter \Delta$
\DisplayProof
&
\Axiom$ \Gamma \fCenter \Delta, A(t) $
\RightLabel{\RightR{\lexists}}
\UnaryInf$ \Gamma \fCenter \Delta, \lexists[x][A(x)]$
\DisplayProof
\end{tabular}
\caption{Rules of \Log{LK}. In $\RightR{\forall}$ and
$\LeftR{\exists}$, the \olup{constant}~$a$ must not occur in the
conclusion sequent. $\Gamma$, $\Delta$, and~$\Pi$ are
sequences of \olup{formula}s.}
\ollabel{tab:G1c}
\end{table}

% ==== OLP file: content/proof-theory/sequent-calculus/rules-mG3i.tex ====
% Part: proof-theory
% Chapter: sequent-calculus
% Section: rules-G3c



\begin{table}
\begin{tabular}{@{}c@{}c@{}}
\multicolumn{2}{@{}c@{}}{Axioms}\\[2ex]
$C, \Gamma \Sequent \Delta, C$ & $\lfalse, \Gamma \Sequent \Delta$
\\[2ex]
\multicolumn{2}{@{}c@{}}{Logical Rules}\\[2ex]
\Axiom$\lnot A, \Gamma \fCenter \Delta, A $
\RightLabel{\LeftR{\lnot}}
\UnaryInf$\lnot A, \Gamma \fCenter \Delta$
\DisplayProof
&
\Axiom$A, \Gamma \fCenter$
\RightLabel{\RightR{\lnot}}
\UnaryInf$ \Gamma \fCenter \Delta, \lnot A $
\DisplayProof
\\[4ex]
\Axiom$ A, B, \Gamma \fCenter \Delta$
\RightLabel{\LeftR{\land}}
\UnaryInf$ A \land B, \Gamma \fCenter \Delta$
\DisplayProof
&
\Axiom$\Gamma \fCenter \Delta, A$
\Axiom$ \Gamma \fCenter \Delta, B$
\RightLabel{\RightR{\land}}
\BinaryInf$ \Gamma \fCenter \Delta, A \land B $
\DisplayProof
\\[4ex]\Axiom$A, \Gamma \fCenter \Delta$
\Axiom$B, \Gamma \fCenter \Delta$
\RightLabel{\LeftR{\lor}}
\BinaryInf$A \lor B, \Gamma \fCenter \Delta$
\DisplayProof
&
\Axiom$\Gamma \fCenter \Delta, A, B$
\RightLabel{\RightR{\lor}}
\UnaryInf$ \Gamma \fCenter \Delta, A \lor B$
\DisplayProof
\\[4ex]
\Axiom$A \lif B, \Gamma \fCenter \Delta, A$
\Axiom$ B, \Gamma \fCenter \Delta$
\RightLabel{\LeftR{\lif}}
\BinaryInf$A \lif B, \Gamma \fCenter \Delta$
\DisplayProof
&
\Axiom$ A, \Gamma \fCenter B$
\RightLabel{\RightR{\lif}}
\UnaryInf$ \Gamma \fCenter \Delta, A \lif B $
\DisplayProof
\\[4ex]
\Axiom$ A(t), \lforall[x][A(x)]\Gamma \fCenter \Delta$
\RightLabel{\LeftR{\lforall}}
\UnaryInf$ \lforall[x][A(x)],\Gamma \fCenter \Delta$
\DisplayProof
&
\Axiom$ \Gamma \fCenter A(a) $
\RightLabel{\RightR{\lforall}}
\UnaryInf$ \Gamma \fCenter \lforall[x][A(x)]$
\DisplayProof
\\[4ex]
\Axiom$ A(a), \Gamma \fCenter \Delta $
\RightLabel{\LeftR{\lexists}}
\UnaryInf$ \lexists[x][A(x)], \Gamma \fCenter \Delta$
\DisplayProof
&
\Axiom$ \Gamma \fCenter \Delta, \lexists[x][A(x)], A(t) $
\RightLabel{\RightR{\lexists}}
\UnaryInf$ \Gamma \fCenter \Delta, \lexists[x][A(x)]$
\DisplayProof
\end{tabular}
\caption{Rules of multi-conclusion \Log{mG3i} for intuitionistic
logic. In $\RightR{\forall}$ and $\LeftR{\exists}$, the
\olup{constant}~$a$ must not occur in the conclusion sequent.
$\Gamma$~and $\Delta$ are multisets of \olup{formula}s. In axioms, the
\olup{formula}~$C$ must be atomic. \Log{mG1m} is
\Log{mG1i} without axioms $\lfalse,
\Gamma \Sequent \Delta$.}
\ollabel{tab:G3c}
\end{table}

% ==== OLP file: content/proof-theory/sequent-calculus/rules-proofs.tex ====
% Part: proof-theory
% Chapter: sequent-calculus
% Section: rules-proofs



\olfileid{pt}{seq}{rp}

\olsection{Rules and \usetoken{P}{proof}}

We begin by giving some definitions and fixing some terminology.

Consider, for example, the two rules in \Log{G3c} that allow \olup{proving}
sequents containing $A \land B$ from sequents containing $A$
and~$B$:
\[
\Axiom$ A, B, \Gamma \fCenter \Delta$
\RightLabel{\LeftR{\land}}
\UnaryInf$ A \land B, \Gamma \fCenter \Delta$
\DisplayProof
\qquad
\Axiom$\Gamma \fCenter \Delta, A$
\Axiom$ \Gamma \fCenter \Delta, B$
\RightLabel{\RightR{\land}}
\BinaryInf$ \Gamma \fCenter \Delta, A \land B $
\DisplayProof
\]
The sequents on top of the line are called the \emph{premises} and the
sequent below the line the \emph{conclusion}. The \olup{formula}s in
$\Gamma$ and $\Delta$ in each rule are called the \emph{context} or
the \emph{side \olup{formula}s}. The \olup{formula} in the conclusion that
isn't part of the context is called the \emph{principal} \olup{formula},
and the \olup{formula}s in the premises not in the context the
\emph{active} (or \emph{auxiliary}) \olup{formula}s.

The quantifier rules are slightly more complicated. E.g., the rules in
\Log{G1c} for $\lforall$ are:
\[\Axiom$ A(t), \Gamma \fCenter \Delta$
\RightLabel{\LeftR{\lforall}}
\UnaryInf$ \lforall[x][A(x)],\Gamma \fCenter \Delta$
\DisplayProof
\qquad
\Axiom$ \Gamma \fCenter \Delta, A(a) $
\RightLabel{\RightR{\lforall}}
\UnaryInf$ \Gamma \fCenter \Delta, \lforall[x][A(x)]$
\DisplayProof
\]
In the \LeftR{\lforall} rule, the active \olup{formula} is~$A(t)$ where
$t$~is a closed term, i.e., a term without variables. Such an
inference is correct---that is, it matches the schematic rule
\LeftR{\lforall}---if there is some \olup{formula}~$A$, some
variable~$x$, and some closed term~$t$ such that the principal
\olup{formula} in the conclusion is $\lforall[x][A]$ and the active
\olup{formula} in the premise is~$\Subst{A}{t}{x}$. The
\RightR{\lforall} rule works the same, except that instead of any
term~$t$ the active \olup{formula} in the premise must
be~$\Subst{A}{a}{x}$, where $a$ is \olup{constant}. This constant is
called the \emph{eigenvariable} of the \RightR{\lforall} inference.
The inference is only correct if the lower sequent does not contain~$a$,
that is, if $a$ does not occur in $\Gamma$, $\Delta$, or~$A$. This is
called the \emph{eigenvariable condition}.

Sequent systems like \Log{G1c} and \Log{G3c} consists of a collection
of such schematically specified rules together with some (also
schematically specified) sequents that count as axioms.
!!^{proof}s in such a system are finite trees of sequents that
are inductively generated from axioms using the rules of inference.
Each leaf is an axiom, and every other sequent follows from the
sequents immediately above it by one of the inference rules of the
system.

\begin{defn}\ollabel{defn:proof}
\emph{!!^{proof}s} in a sequent system are finite trees of
sequents inductively as follows:
\begin{enumerate}
  \item\ollabel{defn:prf:axiom} Every axiom is \olup{proof}.
  \item\ollabel{defn:prf:one-prem} If $\pi_1$ is \olup{proof} ending in
  a sequent~$S_1$, and $S$ is a sequent that follows from~$S_1$ by a
  one-premise rule~$R$, then
  \[
  \AxiomC{}
  \RightLabel{$\pi_1$}
  \DeduceC{$S_1$}
  \RightLabel{$R$}
  \UnaryInfC{$S$}
  \DisplayProof
  \]
  is \olup{proof}.
  \item\ollabel{defn:prf:two-prem} If $\pi_1$ and $\pi_2$ are \olup{proof}s ending in sequents~$S_1$ and~$S_2$, and $S$ is a
  sequent that follows from~$S_1$ and$S_2$ by a two-premise rule~$R$, then
  \[
  \AxiomC{}
  \RightLabel{$\pi_1$}
  \DeduceC{$S_1$}
  \AxiomC{}
  \RightLabel{$\pi_2$}
  \DeduceC{$S_2$}
  \RightLabel{$R$}
  \BinaryInfC{$S$}
  \DisplayProof
  \]
  is \olup{proof}.
\end{enumerate}
\end{defn}

\Olref{defn:proof} defines \olup{proof}s as trees of sequents. We think
of these trees as growing ``upwards.'' The topmost (leaf) nodes of
this tree are axioms and called \emph{initial sequents}. The
bottom-most sequent is called the \emph{end-sequent}. If $\pi$ is
\olup{proof} with end-sequent~$S$ we also write $\pi \Proves S$. If
there is \olup{proof} of a sequent~$S$ we write $\Proves S$, possibly
indicating the system in which the proof exists on the left of the
$\Proves$~symbol, e.g., $\Log{G1c} \Proves A, B \Sequent A \land
B$. Every sequent in \olup{proof}, with the exception of initial
sequents, is the \emph{conclusion} of an inference by some rule~$R$.
The one or two sequents standing immediately above a sequent in
\olup{proof} (i.e., the parents of the node) are called the
\emph{premises} of the inference.

The \olup{proof}s used in the inductive construction of
\olup{proof} are called its \emph{sub-\olup{proof}s}. The
sub-\olup{proof}s ending in the premises of an inference are called
the \emph{immediate sub-\olup{proof}s} of the inference.

It is often necessary to measure the complexity of \olup{proof}s as
a natural number. We could simply count the number of sequents in
\olup{proof}, but often a more useful measure is the \olup{height} of
\olup{proof}: it is the maximal number of sequents between the
end-sequent and an initial sequent. We can also give an inductive
definition.

\begin{defn}
The \emph{\olup{height}}~$\pheight{\pi}$ of \olup{proof}~$\pi$ is
inductively defined as follows.
\begin{enumerate}
  \item If $\pi$ consists of a single axiom, $\pheight{\pi} = 0$.
  \item If $\pi$ ends in an inference with one premise with immediate
  sub-\olup{proof}~$\pi_1$, then $\pheight{\pi} = \pheight{\pi_1} + 1$.
  \item If $\pi$ ends in an inference with two premises with immediate
  sub-\olup{proof}s~$\pi_1$ and~$\pi_2$, then $\pheight{\pi} =
  \max(\pheight{\pi_1}, \pheight{\pi_2}) + 1$.
\end{enumerate}
If a sequent~$S$ has \olup{proof} of height $\le n$, we write $\Proves[n] S$.
\end{defn}

\begin{ex}
In \Log{G3c}, any sequent of the form $C, \Gamma \Sequent \Delta, C$
is an axiom, where $C$ must be atomic. Suppose $D$ and $E$ are
atomic sentences. Then $D, E \Sequent D$ and $D, E \Sequent E$
are instances of the axiom $C, \Gamma \Sequent \Delta, C$. E.g., if
we take $C \ident E$, $\Gamma = \{D\}$ and $\Delta = \emptyset$ in
$C, \Gamma \Sequent \Delta, C$ we get exactly the sequent $D, E
\Sequent E$. (Remember that we are dealing with multi\emph{sets} and
so $D, E \Sequent E$ and $E, D \Sequent D$ are the same
sequent.)

Since $D, E \Sequent D$ is an instance of an axiom of~\Log{G3c}, it
(by itself) counts as \olup{proof} in~\Log{G3c}, and so does $D, E
\Sequent E$ by \olref{defn:proof}\olref{defn:prf:axiom}. Let's
call them $\pi_1$ and~$\pi_2$. By \olref{defn:prf:two-prem}, we can
take these two \olup{proof}s and generate a new one ($\pi_3$) using the
two-premise rule~$\RightR{\land}$:
\[
\Axiom$D, E \fCenter D$
\Axiom$D, E \fCenter E$
\RightLabel{\RightR{\land}}
\BinaryInf$D, E  \fCenter D \land E $
\DisplayProof
\]
From $\pi_3$ in turn we obtain the \olup{proof}~$\pi_4$
\[
\Axiom$D, E \fCenter D$
\Axiom$D, E \fCenter E$
\RightLabel{\RightR{\land}}
\insertBetweenHyps{\hspace{5ex}}
\BinaryInf$D, E  \fCenter D \land E $
\LeftSubproofLabel{$\pi_3$}
\RightLabel{\LeftR{\land}}
\UnaryInf$D \land E  \fCenter D \land E$
\RightSubproofLabel{$\pi_4$}
\DisplayProof
\]
using the one-premise rule~$\LeftR{\land}$ (using
\olref{defn:proof}\olref{defn:prf:one-prem}). The immediate
sub-\olup{proof} of the last inference in~$\pi_4$ is~$\pi_3$. The
immediate sub-\olup{proof}s of the $\RightR{\land}$ inference are $\pi_1$
and~$\pi_2$. The sub-\olup{proof}s of $\pi_4$ are all of $\pi_1$,
$\pi_2$, $\pi_3$, and $\pi_4$ itself. We have $\pheight{\pi_1} =
\pheight{\pi_2} = 0$, $\pheight{\pi_3} = 1$ and $\pheight{\pi_4} = 2$.
We have shown that $\Log{G3c} \Proves[4] D \land E  \fCenter D
\land E$ for any atomic $D$ and~$E$.
\end{ex}

% ==== OLP file: content/proof-theory/sequent-calculus/sequent-calculus.tex ====
% Part: proof-theory
% Chapter: sequent-calculus



\olchapter{pt}{seq}{The Sequent Calculus}

\olimport{introduction}
\olimport{rules-proofs}
\olimport{proof-examples}
\olimport{interpretation-rules}
\olimport{quantifiers}
\olimport{admissible-derivable}
\olimport{invertibility}
\olimport{translations}

\OLEndChapterHook


% ==== OLP file: content/proof-theory/sequent-calculus/translations.tex ====
% Part: proof-theory
% Chapter: sequent-calculus
% Section: translations



\olfileid{pt}{seq}{tr}

\olsection{Translating between \Log{G1c} and \Log{G3c}}

We mentioned that both \Log{G1c} and \Log{G3c} are sound and complete
for classical logic, i.e., they each prove all and only the sequents
true in every \olup{structure}. In particular, they prove the same
sequents. We are now in a position to prove this fact without the
detour through the soundness and completeness theorems, but using
purely proof-theoretically. Such a proof provides
\emph{transformations} of \olup{proof}s in \Log{G1c} into \olup{proof}s in
\Log{G3c} and vice versa.

\begin{prop}
If $\Log{G1c} \Proves S$ then $\Log{G3c} \Proves S$.
\end{prop}

\begin{proof}
  We show that if $\pi$ is \olup{proof} in \Log{G1c}, then there is
  \olup{proof}~$\pi'$ of~$S$ in \Log{G3c}. We do this by induction on $n
  = \pheight{\pi}$.

  If $n = 0$, then $\pi$ is an axiom. The axiom $\lfalse \Sequent
  \quad$ is also an axiom in~\Log{G3c} (take the context $\Gamma,
  \Delta$ to be empty). In \Log{G3c}, axioms $A \Sequent A$ with
  arbitrary, not just atomic~$A$ are allowed. But we proved in
  \olref[ex]{prop:G3c-axioms} that all such sequents are \olup{provable}
  in~\Log{G3c}.
  
  If $n>0$, $\pi$ ends in a rule~$R$. The inductive hypothesis
  guarantees that there are \olup{proof}s of the premises of that rule
  in~\Log{G3c}, i.e., that the immediate sub-\olup{proof}s have
  translations in~\Log{G3c}. If the rule~$R$ is also a rule in
  \Log{G3c}, we apply it to the translated \olup{proof}s of the premises.
  For those rules that are no also rules of~\Log{G3c}, we have already
  shown that they are at least admissible. So, if rule~$R$ was
  $\LeftR{\land}$ or $\RightR{\lor}$, apply
  \olref[adm]{prop:lor-G3-adm}. If it was a weakening rule, apply
  \olref[adm]{prop:weak-G3c-adm}. If it was a contraction rule, apply
  \olref[inv]{prop:G3c-cont-adm}.
\end{proof}

\begin{prop}
If $\Log{G3c} \Proves S$ then $\Log{G1c} \Proves S$.
\end{prop}

\begin{proof}
  We again proceed by induction on the \olup{height} of \olup{proof}~$\pi$
  of $S$. Every axiom of \Log{G3c} can be \olup{prove}d in~\Log{G1c} from
  an axiom plus several applications of the \LeftR{\Weakening} and
  \RightR{\Weakening} rules:
  \[
  \Axiom$C \fCenter C$
  \RightLabel{\LeftR{\Weakening}}
  \Deduce$C, \Gamma \fCenter C$
  \RightLabel{\RightR{\Weakening}}
  \Deduce$C, \Gamma \fCenter \Delta, C$
  \DisplayProof
\qquad
  \Axiom$\lfalse \fCenter$
  \RightLabel{\LeftR{\Weakening}}
  \Deduce$\lfalse, \Gamma \fCenter$
  \RightLabel{\RightR{\Weakening}}
  \Deduce$\lfalse, \Gamma \fCenter \Delta$
  \DisplayProof
  \]
  If $\pi$ ends in a rule~$R$, the inductive hypothesis yields
  \olup{proof}s in \Log{G1c} of the premises, and we can apply the same
  rule~$R$ to those \olup{proof}s. The exceptions are the rules
  \LeftR{\land} and \RightR{\lor}, which differ between \Log{G1c}
  and~\Log{G3c}. The \Log{G3c} versions are derivable in \Log{G1c} as follows:
  \[
  \Axiom$A, B, \Gamma \fCenter \Delta$
  \RightLabel{\LeftR{\land}}
  \UnaryInf$A \land B, B, \Gamma \fCenter \Delta$
  \RightLabel{\LeftR{\land}}
  \UnaryInf$A \land B, A \land B, \Gamma \fCenter \Delta$
  \RightLabel{\LeftR{\Contraction}}
  \UnaryInf$A \land B, \Gamma \fCenter \Delta$
    \DisplayProof
\qquad
  \Axiom$\Gamma \fCenter \Delta, A, B$
  \RightLabel{\RightR{\lor}}
  \UnaryInf$\Gamma \fCenter \Delta, A \lor B, B$
  \RightLabel{\RightR{\lor}}
  \UnaryInf$\Gamma \fCenter \Delta, , A \lor B, A \lor B$
  \RightLabel{\RightR{\Contraction}}
  \UnaryInf$\Gamma \fCenter \Delta, A \lor B$
    \DisplayProof
  \]
\end{proof}

\end{document}
